\section{AIME II 2022}
        \begin{enumerate}
        	\item (\textit{AIME II 2022}) Adults made up $\frac5{12}$ of the crowd of people at a concert. After a bus carrying $50$ more people arrived, adults made up $\frac{11}{25}$ of the people at the concert. Find the minimum number of adults who could have been at the concert after the bus arrived.


 
 \textbf{Answer: }154



\textbf{Solution 1}

Let $x$ be the number of people at the party before the bus arrives. We know that $x\equiv 0\pmod {12}$, as $\frac{5}{12}$ of people at the party before the bus arrives are adults. Similarly, we know that $x + 50 \equiv 0 \pmod{25}$, as $\frac{11}{25}$ of the people at the party are adults after the bus arrives. $x + 50 \equiv 0 \pmod{25}$ can be reduced to $x \equiv 0 \pmod{25}$, and since we are looking for the minimum amount of people, $x$ is $300$. That means there are $350$ people at the party after the bus arrives, and thus there are $350 \cdot \frac{11}{25} = \boxed{154}$ adults at the party.

 





\textbf{Solution 2 (Kind of lame)}

Since at the beginning, adults make up $\frac{5}{12}$ of the concert, the amount of people must be a multiple of 12.

 Call the amount of people in the beginning $x$.Then $x$ must be divisible by 12, in other words: $x$ must be a multiple of 12.Since after 50 more people arrived, adults make up $\frac{11}{25}$ of the concert, $x+50$ is a multiple of 25. This means $x+50$ must be a multiple of 5.

 Notice that if a number is divisible by 5, it must end with a 0 or 5. Since 5 is impossible (obviously, since multiples of 12 end in 2, 4, 6, 8, 0,...), $x$ must end in 0.

 Notice that the multiples of 12 that end in 0 are: 60, 120, 180, etc.. By trying out, you can clearly see that $x=300$ is the minimum number of people at the concert.

 So therefore, after 50 more people arrive, there are $300+50=350$ people at the concert, and the number of adults is $350*\frac{11}{25}=154$. Therefore the answer is $\boxed{154}$. 

 I know this solution is kind of lame, but this is still pretty straightforward. This solution is very similar to the first one, though.

 





\textbf{Solution 3}

Let $a$ be the number of adults before the bus arrived and $x$ be the total number of people at the concert. So, $\frac{a}{x}=\frac{5}{12}$. Solving for $x$ in terms of $a$, $x = \frac{12}{5}a$. After the bus arrives, let's say there are an additional $y$ adults out of the 50 more people who enter the concert. From that, we get $\frac{a+y}{x+50}=\frac{11}{25}$. Replacing $x$ with the value of $a$, the second equation becomes $\frac{a+y}{\frac{12}{5}a+50}=\frac{11}{25}$.

 By cross-multiplying and simplifying, we get that $25(y-22)=\frac{7a}{5}$.

 Observe that we must make sure $y-22$ is positive and divisible by $7$ to have an integer value of $a$. The smallest possible value of $y$ that satisfies this conditions is $29$. Plugging this into the equation, $a = 125$. The question asks for the minimum number of adults that are there after the bus arrives, which is $a+y$. Thus, the answer is simply $125+29=\boxed{154}$.

 





\textbf{Solution 4}

Let $T$ denote the total, $a$ and $b$ be the number of adults and children respectively, $x$ and $y$ be the number of adults and children who arrive by bus respectively.

 
 \[x+y=50\]
 \[\frac{5}{12}T = a\]
 \[\frac{7}{12} = b\]
 \[a+b = T\]
 \[\frac{11}{25}(T+50) = a+x\]
 \[\frac{14}{25} (T+50) = b+y\]

 Consider $\frac{11}{25}(T+50) = a+x$ and substitute $\frac{5}{12}T = a$. We obtain, solving for $x$, $\frac{7}{300}T + 22 = x$.

 Note $y$ is dependent on $x$, $T$, and $b$ but $b$ is dependent on $a$ dependent on $T$. Therefore, the values of $T x$ uniquely determine all other aspects. Here, $T \equiv 0 \pmod{300}$, so minimality occurs $T = 300$, in which $x = 29$, $a = 125$, and $a+x = \boxed{154}$.

 





\textbf{Solution 5}

Let $x$ be the original total number of people, and $a$ be the original number of adults. Let $b$ be the number of adults that arrive on the bus. We can say that $a$ = $\frac{5}{12}x$. This means that for $x$ to be an integer, it must be a multiple of $12$. Let $x$ = $12k$. We can now write a proportion and minimize it for $k$. 

 
 \[\frac{5k+b}{12k+50} = \frac{11}{25}\]

 We solve and simplify this equation to get that $7k$ = $25b - 550$. We want to minimize $k$, and we know that the right side of the equation must be divisible by $7$. We test values of $b$ starting at $23$, we then see that $29$ works. When we plug $29$ into the expression and solve to see that $k$ is $25$. We find the number of adults at the party by substituting $25$ into the expression $5k + 29$ to see that the answer is $\boxed{154}$.

 




	\item (\textit{AIME II 2022}) Azar, Carl, Jon, and Sergey are the four players left in a singles tennis tournament. They are randomly assigned opponents in the semifinal matches, and the winners of those matches play each other in the final match to determine the winner of the tournament. When Azar plays Carl, Azar will win the match with probability $\frac23$. When either Azar or Carl plays either Jon or Sergey, Azar or Carl will win the match with probability $\frac34$. Assume that outcomes of different matches are independent. The probability that Carl will win the tournament is $\frac{p}{q}$, where $p$ and $q$ are relatively prime positive integers. Find $p+q$.


 
 \textbf{Answer: }125



\textbf{Solution}

Let $A$ be Azar, $C$ be Carl, $J$ be Jon, and $S$ be Sergey. The $4$ circles represent the $4$ players, and the arrow is from the winner to the loser with the winning probability as the label.

 
                    \begin{figure}[H]
                        \centering
                        \includegraphics[width=0.5\textwidth]{Solution/AIME_2022_II_2_Solution_Figure_1.png}
                    \end{figure}
                    

 This problem can be solved by using $2$ cases.

 $\textbf{Case 1:}$ $C$'s opponent for the semifinal is $A$

 The probability $C$'s opponent is $A$ is $\frac13$. Therefore the probability $C$ wins the semifinal in this case is $\frac13 \cdot \frac13$. The other semifinal game is played between $J$ and $S$, it doesn't matter who wins because $C$ has the same probability of winning either one. The probability of $C$ winning in the final is $\frac34$, so the probability of $C$ winning the tournament in case 1 is $\frac13 \cdot \frac13 \cdot \frac34$

 
$\textbf{Case 2:}$ $C$'s opponent for the semifinal is $J$ or $S$

 It doesn't matter if $C$'s opponent is $J$ or $S$ because $C$ has the same probability of winning either one. The probability $C$'s opponent is $J$ or $S$ is $\frac23$. Therefore the probability $C$ wins the semifinal in this case is $\frac23 \cdot \frac34$. The other semifinal game is played between $A$ and $J$ or $S$. In this case it matters who wins in the other semifinal game because the probability of $C$ winning $A$ and $J$ or $S$ is different. 

 $\textbf{Case 2.1:}$ $C$'s opponent for the final is $A$

 For this to happen, $A$ must have won $J$ or $S$ in the semifinal, the probability is $\frac34$. Therefore, the probability that $C$ won $A$ in the final is $\frac34 \cdot \frac13$.

 $\textbf{Case 2.2:}$ $C$'s opponent for the final is $J$ or $S$

 For this to happen, $J$ or $S$ must have won $A$ in the semifinal, the probability is $\frac14$. Therefore, the probability that $C$ won $J$ or $S$ in the final is $\frac14 \cdot \frac34$.

 In Case 2 the probability of $C$ winning the tournament is $\frac23 \cdot \frac34 \cdot (\frac34 \cdot \frac13 + \frac14 \cdot \frac34)$

 Adding case 1 and case 2 together we get $\frac13 \cdot \frac13 \cdot \frac34 + \frac23 \cdot \frac34 \cdot (\frac34 \cdot \frac13 + \frac14 \cdot \frac34) = \frac{29}{96},$ so the answer is $29 + 96 = \boxed{\textbf{125}}$.

 




	\item (\textit{AIME II 2022}) A right square pyramid with volume $54$ has a base with side length $6.$ The five vertices of the pyramid all lie on a sphere with radius $\frac mn$, where $m$ and $n$ are relatively prime positive integers. Find $m+n$.


 
 \textbf{Answer: }021



\textbf{Solution 1}


                    \begin{figure}[H]
                        \centering
                        \includegraphics[width=0.5\textwidth]{Solution/AIME_2022_II_3_Solution_Figure_1.png}
                    \end{figure}
                    Although I can't draw the exact picture of this problem, but it is quite easy to imagine that four vertices of the base of this pyramid is on a circle (Radius $\frac{6}{\sqrt{2}} = 3\sqrt{2}$). Since all five vertices are on the sphere, the distances of the spherical center and the vertices are the same: $l$. Because of the symmetrical property of the pyramid,we can imagine that the line of the apex and the (sphere's) center will intersect the square at the (base's) center.

 Since the volume is $54 = \frac{1}{3} \cdot S \cdot h = \frac{1}{3} \cdot 6^2 \cdot h$, where $h=\frac{9}{2}$ is the height of this pyramid, we have: $l^2=\left(\frac{9}{2}-l\right)^2+\left(3\sqrt{2}\right)^2$ according to the Pythagorean theorem.

 Solve this equation will give us $l = \frac{17}{4},$ therefore $m+n=\boxed{021}.$

 





\textbf{Solution 2}

To start, we find the height of the pyramid. By the volume of a pyramid formula, we have 
 \[\frac13 \cdot 6^2 \cdot h=54 \implies h=\frac92.\] Next, let us find the length of the non-base sides of the pyramid. By the Pythagorean Theorem, noting that the distance from one vertex of the base to the center of the base is $\frac12 \cdot 6\sqrt2=3\sqrt2$, we have 
 \[x=\sqrt{\left(\frac92\right)^2+(3\sqrt2)^2}=\sqrt{\frac{153}4}=\frac{3\sqrt{17}}2.\] Taking the cross section of the pyramid and transforming the problem into $2$-d, it suffices to find the radius of the circumcircle of a triangle of side lengths $\frac{3\sqrt{17}}2$, $\frac{3\sqrt{17}}2$, $6\sqrt2$. This turns out to be easy by the formula $R=\frac{abc}{4A}$, and through computing this value (the work has been left out) we find that $R=\frac{17}4$, so our answer is $\boxed{\textbf{021}}$.

 





\textbf{Solution 3 (Coord bash)}

By the volume of a pyramid formula, we have that the height of the pyramid is $\frac{9}{2}$. Since the base is a square with side length 6, the simplest way to place it in the coordinate plane is to put the center of the square at the origin and let the base be on the $xy$ plane. Then, the vertices of the base would be $(3,3,0), (3,-3,0), (-3,3,0), (-3,-3,0)$ in some order. Also, let the vertex be  $(0,0,\frac{9}{2})$. Recall that the formula for a sphere is $(x-a)^2+(y-b)^2+(z-c)^2=r^2$ where the center is $(a,b,c)$ and the radius is $r$. Symmetry gives that $a=b=0$. Plug in $(3,3,0)$ and $(0,0,\frac{9}{2})$ and you get the system of equation 

 $18+c^2=r^2$

 $(\frac{9}{2}-c)^2=r^2$

 Solving gives $c=1/4$ and $r=17/4$, so our answer is $17+4=\boxed{021}$.~\href{https://artofproblemsolving.com/wiki/index.php?title=Ddk001}{Ddk001}

 





\textbf{Solution 4 (Power of a Point)}

We know that the volume of a square pyramid is $\frac{1}{3}\cdot{s^2}\cdot{h}$. The volume of the pyramid is $54$ and the square's side length is $6$. Plugging the information back into the formula, we get

 
 \[\frac{1}{3} \cdot 6^2 \cdot h=54 \implies h=\frac{9}{2}\].

 By Pythagorean theorem, we can say the diagonal of the square is 

 
 \[\sqrt{6^2+6^2}=6\sqrt{2}\]. 

 If we draw a line going through the poles of the sphere, we can see that the line perpendicularly bisects the square's diagonal. If we take the cross-section of the sphere making the problem $2$-d, and call the radius of the sphere $r$, by the power of point, we can say

 
 \[\frac{9}{2}\cdot(2r-\frac{9}{2})=3\sqrt2\cdot{3\sqrt2}\]

 solving this equation gives you

 
 \[r=\frac{17}{4}\]

 so the answer is $17+4=\boxed{021}$. ~\href{https://artofproblemsolving.com/wiki/index.php?title=User:Lentarot}{User:Lentarot}

 




	\item (\textit{AIME II 2022}) There is a positive real number $x$ not equal to either $\frac{1}{20}$ or $\frac{1}{2}$ such that
 \[\log_{20x} (22x)=\log_{2x} (202x).\]The value $\log_{20x} (22x)$ can be written as $\log_{10} (\frac{m}{n})$, where $m$ and $n$ are relatively prime positive integers. Find $m+n$.


 
 \textbf{Answer: }112



\textbf{Solution 1}

Define $a$ to be $\log_{20x} (22x) = \log_{2x} (202x)$, what we are looking for. Then, by the definition of the logarithm,
 \[\begin{cases}  (20x)^{a} &= 22x \\ (2x)^{a} &= 202x.  \end{cases}\]Dividing the first equation by the second equation gives us $10^a = \frac{11}{101}$, so by the definition of logs, $a = \log_{10} \frac{11}{101}$. This is what the problem asked for, so the fraction $\frac{11}{101}$ gives us $m+n = \boxed{112}$.

 





\textbf{Solution 2}

We could assume a variable $v$ which equals to both $\log_{20x} (22x)$ and $\log_{2x} (202x)$.

 So that $(20x)^v=22x \textcircled{1}$and $(2x)^v=202x \textcircled{2}$

 Express $\textcircled{1}$ as: $(20x)^v=(2x \cdot 10)^v=(2x)^v \cdot \left(10^v\right)=22x \textcircled{3}$

 Substitute $\textcircled{{2}}$ to $\textcircled{3}$: $202x \cdot (10^v)=22x$

 Thus, $v=\log_{10} \left(\frac{22x}{202x}\right)= \log_{10} \left(\frac{11}{101}\right)$, where $m=11$ and $n=101$.

 Therefore, $m+n = \boxed{112}$.

 





\textbf{Solution 3}

We have
 \begin{align*} \log_{20x} (22x) & = \frac{\log_k 22x}{\log_k 20x} \\ & = \frac{\log_k x + \log_k 22}{\log_k x + \log_k 20} . \end{align*}


 We have
 \begin{align*} \log_{2x} (202x) & = \frac{\log_k 202x}{\log_k 2x} \\ & = \frac{\log_k x + \log_k 202 }{\log_k x + \log_k 2} . \end{align*}


 Because $\log_{20x} (22x)=\log_{2x} (202x)$, we get
 \[ \frac{\log_k x + \log_k 22}{\log_k x + \log_k 20} = \frac{\log_k x + \log_k 202 }{\log_k x + \log_k 2} . \]

 We denote this common value as $\lambda$.

 By solving the equality $\frac{\log_k x + \log_k 22}{\log_k x + \log_k 20} = \lambda$, we get $\log_k x = \frac{\log_k 22 - \lambda \log_k 20}{\lambda - 1}$.

 By solving the equality $\frac{\log_k x + \log_k 202 }{\log_k x + \log_k 2} = \lambda$, we get $\log_k x = \frac{\log_k 202 - \lambda \log_k 2}{\lambda - 1}$.

 By equating these two equations, we get
 \[ \frac{\log_k 22 - \lambda \log_k 20}{\lambda - 1} = \frac{\log_k 202 - \lambda \log_k 2}{\lambda - 1} . \]

 Therefore,
 \begin{align*} \log_{20x} (22x) & = \lambda \\ & = \frac{\log_k 22 - \log_k 202}{\log_k 20 - \log_k 2} \\ & = \frac{\log_k \frac{11}{101}}{\log_k 10} \\ & = \log_{10} \frac{11}{101} . \end{align*}


 Therefore, the answer is $11 + 101 = \boxed{\textbf{112}}$.

 





\textbf{Solution 4 (Solution 1 with more reasoning)}

Let $a$ be the exponent such that $(20x)^a = 22x$ and $(2x)^a = 202x$. Dividing, we get
 \begin{align*} \dfrac{(20x)^a}{(2x)^a} &= \dfrac{22x}{202x}. \\ \left(\dfrac{20x}{2x}\right)^a &= \dfrac{22x}{202x}. \\ 10^a &= \dfrac{11}{101}. \\ \end{align*}
Thus, we see that $\log_{10} \left(\dfrac{11}{101}\right) = a = \log_{20x} 22x$, so the answer is $11 + 101 = \boxed{112}$.

 





\textbf{Solution 5}

By the change of base rule, we have $\frac{\log 22x}{\log 20x}=\frac{\log 202x}{\log 2x}$, or $\frac{\log 22 +\log x}{\log 20 +\log x}=\frac{\log 202 +\log x}{\log 2 +\log x}=k$. We also know that if $a/b=c/d$, then this also equals $\frac{a-c}{b-d}$. We use this identity and find that $k=\frac{\log 202 -\log 22}{\log 2 -\log 20}=-\log\frac{202}{22}=\log\frac{11}{101}$. The requested sum is $11+101=\boxed{112}.$

 





\textbf{Solution 6}

By change of base formula,
 \[\frac{\log_{2x} 22x}{\log_{2x} 20x} = \frac{{\log_{2x} 11} + 1}{{\log_{2x} 10} + 1} = {\log_{2x} 101} + 1\]
 \[\log_{2x} 11 + 1 = (\log_{2x} 10)(\log_{2x} 101) + \log_{2x} 10 + 1\]
 \[\frac{\log_{2x} \frac{11}{1010}}{\log_{2x} 10} = \log_{2x} 101\]
 \[\log_{10} {\frac{11}{1010}} = \log_{2x} 101\]
 \[\log_{10} {\frac{11}{1010}} + 1 = \log_{2x} 101 + 1 = \log_{2x} 202x = \log_{20x} {22x}\]Thus,
 \[\log_{20x} 22x = \log_{10} \left( \frac{11}{1010} \times 10 \right) = \log_{10} \frac{11}{101}\]The requested answer is $11 + 101 = \boxed{112}$.

 




	\item (\textit{AIME II 2022}) Twenty distinct points are marked on a circle and labeled $1$ through $20$ in clockwise order. A line segment is drawn between every pair of points whose labels differ by a prime number. Find the number of triangles formed whose vertices are among the original $20$ points.


 
 \textbf{Answer: }072



\textbf{Solution 1}

Let $a$, $b$, and $c$ be the vertex of a triangle that satisfies this problem, where $a > b > c$.
 \[a - b = p_1\]
 \[b - c = p_2\]
 \[a - c = p_3\]

 $p_3 = a - c = a - b + b - c = p_1 + p_2$. Because $p_3$ is the sum of two primes, $p_1$ and $p_2$, $p_1$ or $p_2$ must be $2$. Let $p_1 = 2$, then $p_3 = p_2 + 2$. There are only $8$ primes less than $20$: $2, 3, 5, 7, 11, 13, 17, 19$. Only $3, 5, 11, 17$ plus $2$ equals another prime. $p_2 \in \{ 3, 5, 11, 17 \}$.

 Once $a$ is determined, $a = b+2$ and $b = c + p_2$. There are $18$ values of $a$ where $b+2 \le 20$, and $4$ values of $p_2$. Therefore the answer is $18 \cdot 4 = \boxed{\textbf{072}}$

 \textbf{Note: This solution seems incorrect.} Although the answer is correct, solution 2 below is a more accurate way to approach this problem. I agree, I don't get how $a + 2 \leq 20$.

 





\textbf{Solution 2}

As above, we must deduce that the sum of two primes must be equal to the third prime. Then, we can finish the solution using casework.If the primes are $2,3,5$, then the smallest number can range between $1$ and $15$.If the primes are $2,5,7$, then the smallest number can range between $1$ and $13$.If the primes are $2,11,13$, then the smallest number can range between $1$ and $7$. If the primes are $2,17,19$, then the smallest number can only be $1$.

 Adding all cases gets $15+13+7+1=36$. However, due to the commutative property, we must multiply this by 2. For example, in the $2,17,19$ case the numbers can be $1,3,20$ or $1,18,20$. Therefore the answer is $36\cdot2=\boxed{072}$.

 Note about solution 1: I don't think that works, because if for example there are 21 points on the circle, your solution would yield $19\cdot4=76$, but there would be $8$ more solutions than if there are $20$ points. This is because the upper bound for each case increases by $1$, but commutative property doubles it to be $4$.

 




	\item (\textit{AIME II 2022}) Let $x_1\leq x_2\leq \cdots\leq x_{100}$ be real numbers such that $|x_1| + |x_2| + \cdots + |x_{100}| = 1$ and $x_1 + x_2 + \cdots + x_{100} = 0$. Among all such $100$-tuples of numbers, the greatest value that $x_{76} - x_{16}$ can achieve is $\frac mn$, where $m$ and $n$ are relatively prime positive integers. Find $m+n$.


 
 \textbf{Answer: }841



\textbf{Solution 1}

To find the greatest value of $x_{76} - x_{16}$, $x_{76}$ must be as large as possible, and $x_{16}$ must be as small as possible. If $x_{76}$ is as large as possible, $x_{76} = x_{77} = x_{78} = \dots = x_{100} > 0$. If $x_{16}$ is as small as possible, $x_{16} = x_{15} = x_{14} = \dots = x_{1} < 0$. The other numbers between $x_{16}$ and $x_{76}$ equal to $0$. Let $a = x_{76}$, $b = x_{16}$. Substituting $a$ and $b$ into $|x_1| + |x_2| + \cdots + |x_{100}| = 1$ and $x_1 + x_2 + \cdots + x_{100} = 0$ we get:
 \[25a - 16b = 1\]
 \[25a + 16b = 0\]$a = \frac{1}{50}$, $b = -\frac{1}{32}$

 $x_{76} - x_{16} = a - b = \frac{1}{50} + \frac{1}{32} = \frac{41}{800}$. $m+n = \boxed{\textbf{841}}$

 





\textbf{Solution 2}

Define $s_N$ to be the sum of all the negatives, and $s_P$ to be the sum of all the positives.

 Since the sum of the absolute values of all the numbers is $1$, $|s_N|+|s_P|=1$.

 Since the sum of all the numbers is $0$, $s_N=-s_P\implies |s_N|=|s_P|$.

 Therefore, $|s_N|=|s_P|=\frac 12$, so $s_N=-\frac 12$ and $s_P=\frac 12$ since $s_N$ is negative and $s_P$ is positive.

 To maximize $x_{76}-x_{16}$, we need to make $x_{16}$ as small of a negative as possible, and $x_{76}$ as large of a positive as possible.

 Note that $x_{76}+x_{77}+\cdots+x_{100}=\frac 12$ is greater than or equal to $25x_{76}$ because the numbers are in increasing order.

 Similarly, $x_{1}+x_{2}+\cdots+x_{16}=-\frac 12$ is less than or equal to $16x_{16}$.

 So we now know that $\frac 1{50}$ is the best we can do for $x_{76}$, and $-\frac 1{32}$ is the least we can do for $x_{16}$.

 Finally, the maximum value of $x_{76}-x_{16}=\frac 1{50}+\frac 1{32}=\frac{41}{800}$, so the answer is $\boxed{841}$.

 (Indeed, we can easily show that $x_1=x_2=\cdots=x_{16}=-\frac 1{32}$, $x_{17}=x_{18}=\cdots=x_{75}=0$, and $x_{76}=x_{77}=\cdots=x_{100}=\frac 1{50}$ works.)

 





\textbf{Solution 3}

Because the absolute value sum of all the numbers is $1$, and the normal sum of all the numbers is $0$, the positive numbers must add to $\frac12$ and negative ones must add to $-\dfrac12$. To maximize $x_{76} - x_{16}$, we must make $x_{76}$ as big as possible and $x_{16}$ as small as possible. We can do this by making $x_1 + x_2 + x_3 \dots x_{16} = -\dfrac{1}{2}$, where $x_1 = x_2 = x_3 = \dots = x_{16}$ (because that makes $x_{16}$ the smallest possible value), and $x_{76} + x_{77} + x_{78} + \dots + x_{100} = \dfrac{1}{2}$, where similarly $x_{76} = x_{77} = \dots = x_{100}$ (because it makes $x_{76}$ its biggest possible value.) That means $16(x_{16}) = -\dfrac{1}{2}$, and $25(x_{76}) = \dfrac{1}{2}$. $x_{16} = -\dfrac{1}{32}$ and $x_{76} = \dfrac{1}{50}$, and subtracting them $\dfrac{1}{50} - \left( -\dfrac{1}{32}\right) = \dfrac{41}{800}$. $41 + 800 = 841$.

 





\textbf{Solution 4}

All positive terms sum to $\frac{1}{2}$ and all negative to $\frac{-1}{2}$.

 \textbf{Lemma 1:} $x_{16} \geq \frac{-1}{32}$.

 \textbf{Proof of Lemma 1:} We require $x_{16}$ minimal and $x_{76}$ as maximal, which implies because the sequence is increasing, we must set all terms between $x_{17}$ and $x_{75}$ as a constant $c$. It is obvious this $c = 0$. Assume it wasn't, and then if $c$ is positive, we have to subtract some $nc$ from $0.5$ and same cases appear for $c$ negative or a combination of both. It follows that $x_{16} \geq -\frac{\frac{1}{2}}{16} = \frac{-1}{32}$ as desired.

 \textbf{Lemma 2:} $x_{76} \leq \frac{1}{50}$

 For space management, the proof of Lemma 2 is similar to that of Lemma 1.

 Then, set all $x_i$ for $1 \leq i \leq 16$ equal and $x_j$ for $76 \leq j \leq 100$ equal to $-\frac{1}{32}$ and $\frac{1}{50}$ respectively, The maximal difference is then $\frac{41}{800} \implies \boxed{841} \quad \square$.

 




	\item (\textit{AIME II 2022}) A circle with radius $6$ is externally tangent to a circle with radius $24$. Find the area of the triangular region bounded by the three common tangent lines of these two circles.


 
 \textbf{Answer: }192



\textbf{Solution 1}


                    \begin{figure}[H]
                        \centering
                        \includegraphics[width=0.5\textwidth]{Solution/AIME_2022_II_7_Solution_Figure_1.png}
                    \end{figure}
                    

 $r_1 = O_1A = 24$, $r_2 = O_2B = 6$, $AG = BO_2 = r_2 = 6$, $O_1G = r_1 - r_2 = 24 - 6 = 18$, $O_1O_2 = r_1 + r_2 = 30$

 $\triangle O_2BD \sim \triangle O_1GO_2$, $\frac{O_2D}{O_1O_2} = \frac{BO_2}{GO_1}$,  $\frac{O_2D}{30} = \frac{6}{18}$, $O_2D = 10$

 $CD = O_2D + r_2 = 10 + 6 = 16$,

 $EF = 2EC = EA + EB = AB = GO_2 = \sqrt{(O_1O_2)^2-(O_1G)^2} = \sqrt{30^2-18^2} = 24$

 $DEF = \frac12 \cdot EF \cdot CD = \frac12 \cdot 24 \cdot 16 = \boxed{\textbf{192}}$

 \textbf{Alternative Finish} Note that $\triangle{O_1 O_2 G} \sim \triangle{O_1 D A}$ by AA similarity. Setting up the ratio $\frac{18}{24}=\frac{24}{AB+8}$ and then substituting $AB$ with our known values will lead us to the same solution.

 \textbf{mathboy282}

 \textbf{Circles of Similitude Method} \( C \) is an internal point of similitude and thus triangle \( O_1AD \sim O_2BD \) from which \( O_1GO_2 \sim O_2BD \). Then one concludes \( CD = 16 \). From there, it is easy to prove that \( ABO_2G \) is a rectangle, and the longer side length \( AB \) is \( 3BD \). One can let the intersection of the circle \( O_2 \) and \( D \) be \( W \), and thus \( WD = 4 \), in which through Pythagorean Theorem \( BD = 8 \) and \( AB = 24 \). The internal point of similitude always bisects both common tangents from an external tangential point \( D \), and thus \( AE = EB \). Definition of tangents show that \( ED = DF = EB + BD = AB/2 + BD = 12 + 8 = 20 \), and because \( CD = 16 \) and \( EC = 12 \) and the given circumstance lies within a right angled triangle, \( EF = 24 \), leading to an answer of \( \frac{1}{2} \times 24 \times 16 = 12 \times 16 = \boxed{192} \).

 As a bit of a side note, you cannot use the nine point circle here as the circumscribed circle of the larger triangle has a nine point circle that does not inscribe the given triangle. In my opinion it was quite smart that MAA accounted for such, otherwise the problem would be trivial.

 





\textbf{Solution 2}

Let the center of the circle with radius $6$ be labeled $A$ and the center of the circle with radius $24$ be labeled $B$. Drop perpendiculars on the same side of line $AB$ from $A$ and $B$ to each of the tangents at points $C$ and $D$, respectively. Then, let line $AB$ intersect the two diagonal tangents at point $P$. Since $\triangle{APC} \sim \triangle{BPD}$, we have 
 \[\frac{AP}{AP+30}=\frac14 \implies AP=10.\] Next, throw everything on a coordinate plane with $A=(0, 0)$ and $B = (30, 0)$. Then, $P = (-10, 0)$, and if $C = (x, y)$, we have 
 \[(x+10)^2+y^2=64,\] 
 \[x^2+y^2=36.\] Combining these and solving, we get $(x, y)=\left(-\frac{18}5, \frac{24}5\right)$. Notice now that $P$, $C$, and the intersections of the lines $x=6$ (the vertical tangent) with the tangent containing these points are collinear, and thus every slope between a pair of points will have the same slope, which in this case is $\frac{-\frac{18}5+10}{\frac{24}5}=\frac34$. Thus, the other two vertices of the desired triangle are $(6, 12)$ and $(6, -12)$. By the Shoelace Formula, the area of a triangle with coordinates $(-10, 0)$, $(6, 12)$, and $(6, -12)$ is 
 \[\frac12|-120-0-72-72+0-120|=\boxed{\textbf{192}}.\] 

 





\textbf{Solution 3}

(Taking diagram names from Solution 1. Also say the line that passes through $O_1$ and is parallel to line EF, call the points of intersection of that line and the circumference of circle $O_1$ points $X$ and $Y$.)

 First notice that $DO_1$ is a straight line because $DXY$ is an isosceles triangle(or you can realize it by symmetry). That means, because $DO_1$ is a straight line, so angle $BDO_2$ = angle $ADO_1,$ triangle $ADO_1$ is similar to triangle $BDO_2$. Also name $DO_2 = x$. By our similar triangles, $\frac{BO_2}{AO_1} = \frac{1}{4} = \frac{x}{x+30}$. Solving we get $x = 10 = DO_2$. Pythagorean Theorem on triangle $DBO_2$ shows $BD = \sqrt{10^2 - 6^2} = 8$. By similar triangles, $DA = 4 \cdot 8 = 32$ which means $AB = DA - DB = 32 - 8 = 24$. Because $BE = CE = AE, AB = 2 \cdot BE = 24$. $BE = 12,$ which means $CE = 12$. $CD = DO_2$(its value found earlier in this solution) + $CO_2$ ($O_2$ 's radius) $= 10 + 6 = 16$. The area of $DEF$ is $\frac{1}{2} \cdot CD \cdot EF = CD \cdot CE$ (because $CE$ is $\frac{1}{2}$ of $EF$) $= 16 \cdot 12 = 192$.

 





\textbf{Solution 4 (similar to solution 1)}


                    \begin{figure}[H]
                        \centering
                        \includegraphics[width=0.5\textwidth]{Solution/AIME_2022_II_7_Solution_Figure_2.png}
                    \end{figure}
                    

 First, we want to find $O_2D$. We know that $\angle O_1AD = \angle O_2BD = 90^{\circ}$, so by AA similarity, $\triangle O_1AD \sim \triangle O_2BD$. We want to find the length of $x$, and using the similar triangles, we write an equation: $\frac{30 + x}{4} = x$. Solving, we get $x=10$. Therefore, $CD = 10 + 6 = 16$. Next, we find that using AA similarity, $\triangle O_2BD \sim \triangle HO_2D \sim \triangle ECD$ and they are 3-4-5 triangles. We can quickly compute $EF = 2EC = 2 \cdot \left( \frac{3}{4} \cdot 16 \right) = 2 \cdot 12 = 24$. Therefore, the area is $\frac{1}{2} \cdot 16 \cdot 24 = \boxed{192}$.

 





\textbf{Solution 5 (Radical Axis)}

The common internal tangent of the two circles intersects the two external tangents at their midpoints as the internal tangent is the radical axis of the two circles. Therefore, the length of the internal tangent is equal to the length of the external tangents which is $24$ by pythagorean theorem. Let $x$ denote the length of the extension of the external tangents such that they intersect. By similar triangles, we get 
 \[\frac{x}{6} = \frac{x+24}{24}\] Therefore, $x = 8$. Our triangle has side lengths $8+12$, $24$, and $8+12$ making it a $20-20-24$ triangle. Dropping the altitude to the side with length $24$, we get the altitude has length $\sqrt{20^2-12^2} = 16$. Therefore, the area of our triangle is $16\cdot 12 = \fbox{192}$.~ Vedoral

 





\textbf{Solution 6 (minimalistic & fast)}


                    \begin{figure}[H]
                        \centering
                        \includegraphics[width=0.5\textwidth]{Solution/AIME_2022_II_7_Solution_Figure_3.png}
                    \end{figure}
                    We are asked for $[PZZ']$.

 


 Obviously, $O_1O_2 = 30$ and $O_1X=6,O_2X=24$

 Observe that since, by definition, $\overline{PY'},\overline{PY}$ are tangents, $\overline{PY'}$ is the reflection of $\overline{PY}$ over $\overline{PO_1}$. So, $Z'$ is the reflection of $Z$ over $\overline{P0_1}$ since that is perpendicular to the line $O_1O_2$ by tangency.

 Therefore, $[PZZ']=2[PZW]$

 let $PO_2=x$

 since $\overline{PX}$ is tangent to $O_1$ at X, and $\overline{PXY}$ is tangent to $O_2$ at $Y$, and $\overline{ZW}$ is tangent to $O_1$, $\angle{PXO_1}=\angle{PYO_2}=\angle{ZWO_1}=90$°$\angle{XPO_1}$ is shared so $\triangle{PXO_1}\sim\triangle{PZW}\sim\triangle{PYO_2}$

 and thus:
 \[\frac{PO_2}{O_2X}=\frac{PO_1}{O_1Y}\implies\frac{x}{6}=\frac{x+30}{24} \implies x =10\implies PX=8\]so $\triangle{PXO_1},\triangle{PZW},\triangle{PYO_2}$ are all 3-4-5 triangles.

 




	\item (\textit{AIME II 2022}) Find the number of positive integers $n \le 600$ whose value can be uniquely determined when the values of $\left\lfloor \frac n4\right\rfloor$, $\left\lfloor\frac n5\right\rfloor$, and $\left\lfloor\frac n6\right\rfloor$ are given, where $\lfloor x \rfloor$ denotes the greatest integer less than or equal to the real number $x$.


 
 \textbf{Answer: }080 or 081 (both were accepted)



\textbf{Solution 1}

We need to find all numbers between $1$ and $600$ inclusive that are multiples of $4$, $5$, and/or $6$ which are also multiples of $4$, $5$, and/or $6$ when $1$ is added to them.

 We begin by noting that the LCM of $4$, $5$, and $6$ is $60$.  We can therefore simplify the problem by finding all such numbers described above between $1$ and $60$ and multiplying the quantity of such numbers by $10$ ($600$/$60$ = $10$).

 After making a simple list of the numbers between $1$ and $60$ and going through it, we see that the numbers meeting this condition are $4$, $5$, $15$, $24$, $35$, $44$, $54$, and $55$.  This gives us $8$ numbers.  $8$ * $10$ = $\boxed{080}$.

 





\textbf{Solution 1.5}

This is Solution 1 with a slick element included.  Solution 1 uses the concept that $60k+l$ is a solution for $n$ if $60k+l$ is a multiple of $3$, $4$, and/or $5$ and $60k+l+1$ is a multiple of $3$, $4$, and/or $5$ for positive integer values of $l$ and essentially any integer value of $k$.  But keeping the same conditions in mind for $k$ and $l$, we can also say that if $60k+l$ is a solution, then $60k-l-1$ is a solution!  Therefore, one doesn't have to go as far as determining the number of values between $1$ and $60$ and then multiplying by $10$.  One only has to determine the number of values between $1$ and $30$ and then multiply by $20$.  The values of $n$ that work between $1$ and $30$ are $4$, $5$, $15$, and $24$.  This gives us $4$ numbers.  $4$ * $20$ = $\boxed{080}$.

 \textbf{Note} Soon after the test was administered, a formal request was made to also accept $\boxed{081}$ as an answer and MAA decided to honor this request. The gist of this request stated that the phrasing of the first part of the question could reasonably be interpreted to mean that one is given the condition to begin with that the integer is less than or equal to $600$.  In this case, if one was told that the values of $\left\lfloor \frac n4\right\rfloor$, $\left\lfloor\frac n5\right\rfloor$, and $\left\lfloor\frac n6\right\rfloor$ were $150$, $120$, and $100$ respectively, then the only possible choice for $n$ would be $600$ as $601$, $602$, and $603$ do not meet the condition as stated in the first part of the problem. If instead the problem asked for the numbers less than $600$ that met the second condition in the problem, the answer would be $\boxed{080}$.  ~burkinafaso ~sethl 

 Here are two interpretations of the problem: 1) If $n$ is a positive integer, how many numbers less than or equal to $600$ are determinable? Answer: $\boxed{080}$. 2) If $n$ is a positive integer and $n \leq 600$, then how many different $n$ are determinable? Answer: $\boxed{081}$. ~epiconan

 





\textbf{Solution 2}

1. For $n$ to be uniquely determined, $n$ AND $n + 1$ both need to be a multiple of $4, 5,$ or $6.$ Since either $n$ or $n + 1$ is odd, we know that either $n$ or $n + 1$ has to be a multiple of $5.$ We can state the following cases: 

 1. $n$ is a multiple of $4$ and $n+1$ is a multiple of $5$

 2. $n$ is a multiple of $6$ and $n+1$ is a multiple of $5$

 3. $n$ is a multiple of $5$ and $n+1$ is a multiple of $4$

 4. $n$ is a multiple of $5$ and $n+1$ is a multiple of $6$

 Solving for each case, we see that there are $30$ possibilities for cases 1 and 3 each, and $20$ possibilities for cases 2 and 4 each. However, we over-counted the cases where 

 1. $n$ is a multiple of $24$ and $n+1$ is a multiple of $5$

 2. $n$ is a multiple of $5$ and $n+1$ is a multiple of $24$

 Each case has $10$ possibilities.

 Adding all the cases and correcting for over-counting, we get $30 + 20 + 30 + 20 - 10 - 10 = \boxed {080}.$

 





\textbf{Solution 2 Supplement}

Here is a detailed solution for Solution 2.

 \begin{verbatim}
$1.$ $4 \mid n$, $5 \mid n+1$$n = 4a$, $n+1 = 4a+1 = 5a - (a - 1)$, $a - 1 = 5k$, $a = 5k+1$, $n = 4(5k+1) = 20k+4$, $20k+4 \le 600$, $20k \le 596$, $k < 30$, $k \in [0, 29]$, 30 integers.
\end{verbatim}\begin{verbatim}
$2.$ $6 \mid n$, $5 \mid n+1$$n = 6a$, $n+1 = 6a+1 = 5a + a + 1$, $a + 1 = 5k$, $a = 5k-1$, $n = 6(5k-1) = 30k-6$, $30k-6 \le 600$, $30k \le 606$, $k \le 20$, $k \in [1, 20]$, 20 integers.
\end{verbatim}\begin{verbatim}
$3.$ $5 \mid n$, $4 \mid n+1$$n = 5a$, $n+1 = 5a+1 = 4a + a + 1$, $a + 1 = 4k$, $a = 4k-1$, $n = 5(4k-1) = 20k-5$, $20k-5 \le 600$, $20k \le 605$, $k \le 30$, $k \in [1, 30]$, 30 integers.
\end{verbatim}\begin{verbatim}
$4.$ $5 \mid n$, $6 \mid n+1$$n = 5a$, $n+1 = 5a+1 = 5(a - 1) + 6$, $a - 1 = 6k$, $a = 6k+1$, $n = 5(6k+1) = 30k+5$, $30k+5 \le 600$, $30k \le 595$, $k < 20$, $k \in [0, 19]$, 20 integers.
\end{verbatim}Over-counted cases:

 \begin{verbatim}
$1.$ $24 \mid n$, $5 \mid n+1$$n = 24a$, $n + 1= 24a + 1 = 20a + 4a + 1$, $4a+1 = 5k$, $k = 2p+1$, $4a = 5(2p + 1) - 1 = 10p + 4$, $n = 24a = 6(10p + 4) = 60p+24$, $60p + 24 \le 600$, $5p + 2 \le 50$, $5p \le 48$, $p < 10$, $p \in [0, 9]$, 10 integers.
\end{verbatim}\begin{verbatim}
$2.$ $5 \mid n$, $24 \mid n+1$$n + 1= 5a + 1 = 24a$, $n = 24a - 1 = 20a + 4a - 1$, $4a-1 = 5k$, $k = 2p+1$, $4a = 5(2p + 1) + 1 = 10p + 6$, $n = 24a = 6(10p + 6) = 60p+36$, $60p + 36 \le 600$, $5p + 3 \le 50$, $5p \le 47$, $p < 10$, $k \in [0, 9]$, 10 integers.
\end{verbatim}$30 + 20 + 30 + 20 - 10 - 10 = \boxed{\textbf{080}}$

 





\textbf{Solution 3}

The problem is the same as asking how many unique sets of values of $\lfloor\frac{n}{4}\rfloor$, $\lfloor\frac{n}{5}\rfloor$, and $\lfloor\frac{n}{6}\rfloor$ can be produced by one and only one value of $n$ for positive integers $n$ less than or equal to 600.

 Seeing that we are dealing with the unique values of the floor function, we ought to examine when it is about to change values, for instance, when $n$ is close to a multiple of 4 in $\lfloor\frac{n}{4}\rfloor$.

 For a particular value of $n$, let $a$, $b$, and $c$ be the original values of $\lfloor\frac{n}{4}\rfloor$, $\lfloor\frac{n}{5}\rfloor$, and $\lfloor\frac{n}{6}\rfloor$, respectively.

 Notice when $n$ $\equiv0\mod4$ and $n$ $\equiv4\mod5$, the value of $\lfloor\frac{n-1}{4}\rfloor$ will be 1 less than the original $a$. The value of $\lfloor\frac{n+1}{5}\rfloor$ will be 1 greater than the original value of $b$.

 More importantly, this means that no other value less than or greater than $n$ will be able to produce the set of original values of $a$, $b$, and $c$, since they make either $a$ or $b$ differ by at least 1.

 
Generalizing, we find that $n$ must satisfy:

 \begin{verbatim}
										$n$ $\equiv0\mod$ $j$
\end{verbatim}\begin{verbatim}
										$n$ $\equiv{k-1}\mod$ $k$
\end{verbatim}Where $j$ and $k$ are pairs of distinct values of 4, 5, and 6.

 Plugging in the values of $j$ and $k$, finding the solutions to the 6 systems of linear congruences, and correcting for the repeated values, we find that there are $\boxed{080}$ solutions of $n$.

 





\textbf{Solution 3 Supplement}

By \href{https://artofproblemsolving.com/wiki/index.php/Chinese\_Remainder\_Theorem}{Chinese Remainder Theorem}, the general solution of systems of $2$ linear congruences is:

 \begin{verbatim}
$n \equiv r_1 (\bmod { \quad m_1})$, $n \equiv r_2 (\bmod { \quad m_2})$, $(m_1, m_2) = 1$Find $k_1$ and $k_2$ such that $k_1 m_1 \equiv 1 (\bmod{ \quad m_2})$, $k_2 m_2 \equiv 1 (\bmod{ \quad m_1})$Then $n \equiv k_2 m_2 r_1 + k_1 m_1 r_2 (\bmod{ \quad m_1 m_2})$
\end{verbatim}$\operatorname{lcm}[4, 5, 6] = 60$, we solve the number of values for $n \le 60$, then multiply by $10$ to get the number of values for $n \le 600$. We are going to solve the following $6$ systems of linear congruences:

 \begin{verbatim}
$\begin{cases} n \equiv 0 \mod{4} \\ n \equiv -1 \mod{5} \end{cases}$$n \equiv 1 \cdot 5 \cdot 0 + 4 \cdot 4 \cdot (-1) \equiv -16 \equiv 4 \mod{20}$, $n \in \{ 4, 24, 44 \}$
\end{verbatim}\begin{verbatim}
$\begin{cases} n \equiv 0 \mod{4} \\ n \equiv -1 \mod{6} \end{cases}$No solution
\end{verbatim}\begin{verbatim}
$\begin{cases} n \equiv 0 \mod{5} \\ n \equiv -1 \mod{4} \end{cases}$$n \equiv 4 \cdot 4 \cdot 0 + 5 \cdot 5 \cdot (-1) \equiv -5 \equiv 15 \mod{20}$, $n \in \{ 15, 35, 55 \}$
\end{verbatim}\begin{verbatim}
$\begin{cases} n \equiv 0 \mod{5} \\ n \equiv -1 \mod{6} \end{cases}$$n \equiv 1 \cdot 6 \cdot 0 + 5 \cdot 5 \cdot (-1) \equiv -25 \equiv 5 \mod{30}$, $n \in \{ 5, 35 \}$
\end{verbatim}\begin{verbatim}
$\begin{cases} n \equiv 0 \mod{6} \\ n \equiv -1 \mod{4} \end{cases}$No solution
\end{verbatim}\begin{verbatim}
$\begin{cases} n \equiv 0 \mod{6} \\ n \equiv -1 \mod{5} \end{cases}$$n \equiv 5 \cdot 5 \cdot 0 + 1 \cdot 6 \cdot (-1) \equiv -6 \equiv 24 \mod{30}$, $n \in \{ 24, 54 \}$
\end{verbatim}$n \in \{ 4, 5, 15, 24, 35, 44, 54, 55 \}$, there are $8$ values for $n \le 60$. For $n \le 600$, the answer is $8 \cdot 10 = \boxed{\textbf{080}}$.

 





\textbf{Solution 4}

Observe that if $1 \le n \le 60$ such that n is a solution to the desired equation, so is $n + 60\cdot m$, where m is an integer, $0 \le m \le 9$.So we only need to consider n from 1 to 60.As shown in Solution 2, there are 4 cases which we will split into 2 main cases:

 
                \begin{itemize}
                    \item Case 1: $4 \mid n$ or $6 \mid n$, $5 \mid (n+1)$\item Case 2: $4 \mid (n+1)$ or $6 \mid (n+1)$, $5 \mid n$
                \end{itemize}


                There are 4 values of n where $1 \le n \le 12$ satisfying $4 \mid n$ or $6 \mid n$.

 I claim that there are 4 values of $n \le 60$ satisfying Case 1. Suppose x is one value of n satisfying $4 \mid n$ or $6 \mid n$, and $n \le 12$.Hence the solutions satisfying $4 \mid n$ or $6 \mid n$, $n \le 60$ are of the form $x + 12m$, so the values of $n + 1$ are $x + 12m + 1 \equiv x + 2m + 1 \equiv 0$ (mod 5), so $2m \equiv 4 + 4x$ (mod 5) and hence the value of m is unique since $0 \le m \le 4$ to satisfy $1 \le n \le 60$ and 2 and 5 are relatively prime.

 A similar approach can be used to show the same for Case 2, that there are 4 values of $n \le 60$.

 Hence our answer is $(4+4)*10 = \fbox{080}$.

 





\textbf{Solution 5}

$n$ and $n + 1$ must be multiples of $4, 5,$ or $6.$ Note that multiples of $4$ and $6$ cannot be consecutive since they're both even.  

 
\textbf{Case 1:} $n \equiv 0 \pmod{4}$ and $n + 1 \equiv 0 \pmod{5}$
 \[n = 4x \implies 4x + 1 \equiv 0 \pmod{5} \implies x \equiv 1\pmod{5}.\] 

 
 \[x = 5a_1 + 1 \implies n = 20a_1 + 4.\]

 $0 \leq a_1 \leq 29 \implies 30$ values for $n$ 

 
\textbf{Case 2:} $n \equiv 0 \pmod{5}$ and $n + 1 \equiv 0 \pmod{4}$
 \[n = 5x \implies 5x + 1 \equiv 3 \pmod{4} \implies x \equiv 3\pmod{4}.\]

 
 \[x = 4a_2 + 3 \implies n = 20a_2 + 15.\]

 $0 \leq a_2 \leq 29 \implies 30$ values for $n$ 

 
\textbf{Case 3:} $n \equiv 0 \pmod{6}$ and $n + 1 \equiv 0 \pmod{5}$
 \[n = 6x \implies 6x \equiv 4 \pmod{5} \implies x \equiv 4\pmod{5}.\]

 
 \[x = 5a_3 + 4 \implies n = 30a_3 + 24.\] 

 $0 \leq a_3 \leq 19 \implies 20$ values for $n$ 

 \textbf{Case 4:} $n \equiv 0 \pmod{5}$ and $n + 1 \equiv 0 \pmod{6}$

 
 \[n = 5x \implies 5x \equiv 5 \pmod{6} \implies x \equiv 1\pmod{6}.\]

 
 \[x = 6a_4 + 1 \implies n = 30a_4 + 5.\]

 $0 \leq a_4 \leq 19 \implies 20$ values for $n$ 

 
\textbf{Case 1 and Case 3 overcount:} 

 $n \equiv 0 \pmod{24}$ and $n + 1 \equiv 0 \pmod{5}$
 \[20a_1 + 4 = 30a_3 + 24 \implies 20a_1 = 30a_3 + 20\] $a_1 = 1,4,7 \dots 28 \implies 10$ overcounted values for these cases. 

 
\textbf{Case 2 and Case 4 overcount:} 

 $n \equiv 0 \pmod{5}$ and $n + 1 \equiv 0 \pmod{24}$
 \[20a_2 + 15 = 30a_4 + 5 \implies 20a_2 = 30a_4 -10\] $a_2 = 1,4,7 \dots 28 \implies 10$ overcounted values for these cases. 

 
The answer is $30 + 30 + 20 + 20 - 10 - 10 = \boxed{80}.$

 





\textbf{Solution 6(Recursive Way)}

Let $n = 60x + y$. Let $\lfloor \frac{n}{4} \rfloor = a$ and define the others analogously. Then we see that $15x + \lfloor \frac{y}{4} \rfloor = a$. We let $y = 60x_1 + y_1$ and we proceed. Eventually, we note that any $y_i$ must be in the range $0$ to $60$ since it can't be greater than $60$. We get the recursive definition of $y_i = 60x_{i + 1} + y_{i + 1}$. Simplifying the series, we would get $15x_1 + 15x_2 + .... + 15x_n = a$. Now note that $\lfloor \frac{n}{4} \rfloor$ has a maximum of $n = 600$ and a minimum of $n = 1$ hence $0 \leq a \leq 150$. Therefore, plugging this in, we have $0 \leq x_1 + ... + x_n \leq 10$. From $y_i = 60x_{i + 1} + y_{i + 1}$, we can deduce $x_{i + 1} = \frac{y_i - y_{i + 1}}{60}$. Plugging this in, we get $0 \leq \frac{y_1 - y_2}{60} + \frac{y_2 - y_3}{60} + ... + \frac{y_{n - 1} - y_n}{60} \leq 10$. This is telescoping series which simplifies to $0 \leq y - y_n \leq 600$. Now this looks almost identical to the original condition which stated $0 \leq n \leq 600$. In fact, we note that any $y_i < 60$ from the floor function condition. Hence, the number of solutions we would have in the range of $0$ to $60$ is proportional to the number of solutions we would have from $0$ to $600$. Since the first set is $60$ units long and the second set is $600$ units long, we note that $\frac{600}{60} = 10$ so when we find the number of solutions between $0$ and $60$, we can just multiply that by $10$ to get our answer.

 The problem reduces to stating:

 How many possible positive integer $n$ less than or equal to $60$ are there such that when given the values of $\lfloor \frac{n}{4} \rfloor, \lfloor \frac{n}{5} \rfloor, \lfloor \frac{n}{6} \rfloor$, we can uniquely find the value of $n$?

 We'll solve this by noting ranges of $4$. Let's first look at $n$ being in the range $[1, 3]$. All the floor values will be $0$ for each of them so no value here is unique. From $[4, 7]$, $n = 4$ will lead values of $1, 0, 0$ respectively. $n = 5$ will lead values of $1, 1, 0$ respectively. And both $n = 6$ and $n = 7$ will lead values of $1, 1, 1$ respectively. So only $n = 4$ and $n = 5$ are unique so we have $2$ cases here. We can perform similar analysis to the range of $[8, 11]$ and find none are unique. We will then do the same for $[12, 15]$. We will see only $n = 15$ gives a unique solution so we have $1$ solution. Nothing will work for $[16, 19]$. Nothing will work for $[20, 23]$. Only $n = 24$ will work for $[24, 27]$. So far the only ones that have worked so far are $n = 4, 5, 15, 24$ out of the first $27$ numbers. We will continue up until the first $30$ numbers then double the ones that have worked so far by symmetry. In the next range, $[28, 31]$ nothing works. Hence, as this range includes the first $30$ numbers, only $4$ have worked so far. If you are in contest, you should be pretty safe assuming the symmetry for the rest of the $30$ numbers. Assuming that is true, we conclude the total numbers that will work is $4 \cdot 2 = 8$. 

 Hence, only $8$ numbers will give a unique solution out of the first $60$ numbers. By our claim using recursion above, we have to multiply this by $10$ to get the actual answer. Hence $8 \cdot 10 = \boxed{80}$.

 \textbf{Note} Important observation is that a multiple of 4 and multiple of 6 cannot be consecutive.

 




	\item (\textit{AIME II 2022}) Let $\ell_A$ and $\ell_B$ be two distinct parallel lines. For positive integers $m$ and $n$, distinct points $A_1, A_2, \allowbreak A_3, \allowbreak \ldots, \allowbreak A_m$ lie on $\ell_A$, and distinct points $B_1, B_2, B_3, \ldots, B_n$ lie on $\ell_B$. Additionally, when segments $\overline{A_iB_j}$ are drawn for all $i=1,2,3,\ldots, m$ and $j=1,\allowbreak 2,\allowbreak 3, \ldots, \allowbreak n$, no point strictly between $\ell_A$ and $\ell_B$ lies on more than two of the segments. Find the number of bounded regions into which this figure divides the plane when $m=7$ and $n=5$. The figure shows that there are 8 regions when $m=3$ and $n=2$.

                    \begin{figure}[H]
                        \centering
                        \includegraphics[width=0.5\textwidth]{Problem/AIME_2022_II_9_Problem_Figure_1.png}
                    \end{figure}
                    


 
 \textbf{Answer: }244



\textbf{Solution 1}

We can use recursion to solve this problem: 

 \textbf{1.} Fix 7 points on $\ell_A$, then put one point $B_1$ on $\ell_B$. Now, introduce a function $f(x)$ that indicates the number of regions created, where x is the number of points on $\ell_B$. For example, $f(1) = 6$ because there are 6 regions.

 \textbf{2.} Now, put the second point $B_2$ on $\ell_B$. Join $A_1~A_7$ and $B_2$ will create $7$ new regions (and we are not going to count them again), and split the existing regions. Let's focus on the spliting process: line segment formed between $B_2$ and $A_1$ intersect lines $\overline{B_1A_2}$, $\overline{B_1A_3}$, ..., $\overline{B_1A_7}$ at $6$ points $\Longrightarrow$ creating $6$ regions (we already count one region at first), then $5$ points $\Longrightarrow$ creating $5$ regions (we already count one region at first), 4 points, etc. So, we have: 
 \[f(2) = f(1) + 7 + (6+5+...+1) = 34.\]

 \textbf{3.} If you still need one step to understand this: $A_1~A_7$ and $B_3$ will still create $7$ new regions. Intersecting 
 \[\overline{A_2B_1}, \overline{A_2B_2};\] 
 \[\overline{A_3B_1}, \overline{A_3B_2};\] 
 \[...\] 
 \[\overline{A_7B_1}, \overline{A_7B_2}\] at $12$ points, creating $12$ regions, etc. Thus, we have: 
 \[f(3) = f(2)+7+(12+10+8+...+2)=34+7+6\cdot 7=83.\]

 Yes, you might already notice that: 
 \[f(n+1) = f(n)+7+(6+5+...+1)\cdot n = f(n) + 7 + 21n.\]

 \textbf{5.} Finally, we have $f(4) = 153$, and $f(5)=244$. Therefore, the answer is $\boxed{244}$.

 Note: we could deduce a general formula of this recursion: $f(n+1)=f(n)+N_a+\frac{n\cdot (N_a) \cdot (N_a-1)}{2}$, where $N_a$ is the number of points on $\ell_A$

 





\textbf{Solution 2}

We want to derive a general function $f(m,n)$ that indicates the number of bounded regions. Observing symmetry, we know this is a symmetric function about $m$ and $n$. Now let's focus on $f(m+1, n)-f(m, n)$, which is the difference caused by adding one point to the existing $m$ points of line $\ell_A$. This new point, call it \#m, when connected to point \#1 on $\ell_B$, crosses $m*(n-1)$ lines, thus making additional $m*(n-1)+1$ bounded regions; when connected to point \#2 on $\ell_B$, it crosses $m*(n-2)$ lines, thus making additional $m*(n-2)+1$ bounded regions; etc. By simple algebra/recursion methods, we see

 $f(m+1, n)-f(m, n)=m*\frac{n(n-1)}{2} +n$

 Notice $f(1,n)=n-1$. Not very difficult to figure out:

 $f(m, n)=\frac{m(m-1)n(n-1)}{4} +mn-1$

 The fact that $f(3,2)=8$ makes us more confident about the formula. Now plug in $m=5, n=7$, we get the final answer of $\boxed{244}$.

 





\textbf{Solution 3}


                    \begin{figure}[H]
                        \centering
                        \includegraphics[width=0.5\textwidth]{Solution/AIME_2022_II_9_Solution_Figure_1.png}
                    \end{figure}
                    Let some number of segments be constructed. We construct a new segment. We start from the straight line $l_B.$ WLOG from point $B_3.$  Segment will cross several existing segments (points $A,B,C,...$) and enter one of the points of the line  $l_A (A_1).$

 Each of these points adds exactly <i>\textbf{ 1 new bounded region}</i> (yellow bounded regions). 

 The exception is the only first segment $(A_1 B_1),$ which does not create any bounded region. Thus, the number of bounded regions is $1$ less than the number of points of intersection of the segments plus the number of points of arrival of the segments to $l_A.$

 Each point of intersection of two segments is determined <i>\textbf{ uniquely}</i> by the choice of pairs of points on each line. 

 The number of such pairs is $\dbinom{n}{2} \cdot \dbinom{m}{2}.$

 Exactly one segment comes to each of the $n$ points of the line $l_A$ from each of the $m$ points of the line $l_B.$ The total number of arrivals is equal to $mn.$Hence, the total number of bounded regions is $N = \dbinom{n}{2} \cdot \dbinom{m}{2} + mn – 1.$

 We plug in $m=5, n=7$, we get the final answer of $\boxed{244}$.

 \textbf{vladimir.shelomovskii@gmail.com, vvsss}

 





\textbf{Solution 4 (Recursion and Complementary Counting)}

When a new point is added to a line, the number of newly bounded regions it creates with each line segment will be one more than the number of intersection points the line makes with other lines.

 Case 1: If a new point $P$ is added to the right on a line when both lines have an equal amount of points.

 WLOG, let the point be on line $\ell_A$. We consider the complement, where new lines don't intersect other line segments. Simply observing, we see that the only line segments that don't intersect with the new lines are lines attached to some point that a new line does not pass through. If we look at a series of points on line $\ell_B$ from left to right and a line connects $P$ to an arbitrary point, then the lines formed with that point and with remaining points on the left of that point never intersect with the line with $P$. Let there be $s$ points on lines $\ell_A$ and $\ell_B$ before $P$ was added. For each of the $s$ points on $\ell_B$, we subtract the total number of lines formed, which is $s^2$, not counting $P$. Considering all possible points on $\ell_B$, we get $(s^2-s)+(s^2-2s)\cdots(s^2-s^2)$ total intersections. However, for each of the lines, there is one more bounded region than number of intersections, so we add $s$. Simplifying, we get $s^3-s\sum_{i=1}^{s}{i}+s\Longrightarrow s(s^2-\sum_{i=1}^{s}{i}+1)$. Note that this is only a recursion formula to find the number of new regions added for a new point $P$ added to $\ell_A$.

 Case 2:  If a new point $P$ is added to the right of a line that has one less point than the other line.

 Continuing on case one, let this point $P$ be on line $\ell_B$. With similar reasoning, we see that the idea remains the same, except $s+1$ lines are formed with $P$ instead of just $s$ lines. Once again, each line from $P$ to a point on line $\ell_A$ creates $s$ non-intersecting lines for that point and each point to its left. Subtracting from $s(s+1)$ lines and considering all possible lines created by $P$, we get $(s(s+1)-s)+(s(s+1)-2s)\cdots(s(s+1)-s(s+1)$ intersections. However, the number of newly bounded regions is the number of intersections plus the number of points on line $\ell_A$. Simplying, we get $s(s+1)^2-s\sum_{i=1}^{s+1}{i}+(s+1)$ newly bounded regions.

 For the base case $s=2$ for both lines, there are $4$ bounded regions. Next, we plug in $s=2,3,4$ for both formulas and plug $s=5$ for the first formula to find the number of regions when $m=6$ and $n=5$. Notice that adding a final point on $\ell_A$ is a variation of our Case 1. The only difference is for each of the $s$ lines formed by $P$, there are $s+1$ points that can form a non-intersecting line. Therefore, we are subtracting a factor of $s+1$ lines instead of $s$ lines from a total of $s(s+1)$ lines. However, the number of lines formed by $P$ remains the same so we still add $s$ at the end when considering intersection points. Thus, the recursive equation becomes $(s(s+1)-(s+1))+(s(s+1)-2(s+1))\cdots(s(s+1)-s(s+1))+s\Longrightarrow s^2(s+1)-(s+1)\sum_{i=1}^{s}{i}+s$. Plugging $s=5$ into this formula and adding the values we obtained from the other formulas, the final answer is $4+4+9+12+22+28+45+55+65=\boxed{244}$.

 





\textbf{Solution 5 (Euler's Graph Formula, similar to solution 3)}

We know the by Euler's Formula for planar graphs that $F-E+V=2$, where $F$ is the number of bounded faces, plus the outer region, $E$ is the number of edges, and $V$ is the number of vertices. Temporarily disregarding the intersections between the lines, we can easily calculate that:

 $V_{i}=7+5=12$

 $E_{i}=6+4+7\cdot5=45$

 However, the resulting graph is not planar, as the edges clearly intersect. To account for this, we must turn all intersection points into vertices, and update our values accordingly.

 Observe that each intersection point can be mapped to two points on either line, and analogously, two points on either line can be mapped to one intersection point, uniquely. Thus, to count intersection points, we simply calculate:

 ${7 \choose 2}{5 \choose 2} = 210$

 And thus,

 $V=V_{i}+210=222$

 We must also account for the edges. Observe that each intersection point turns the two edges that make it into four, that is, each intersection point adds $2$ to the number of edges. Therefore,

 $E=E_{i}+2\cdot210=465$

 Plugging these into Euler's Formula we get:$F-E+V=2$

 $F-465+222=2$

 $F=245$

 Disregarding the outer region, we conclude that our answer is $F-1=245-1=\boxed {244}$

 


 




	\item (\textit{AIME II 2022}) Find the remainder when
 \[\binom{\binom{3}{2}}{2} + \binom{\binom{4}{2}}{2} + \dots +  \binom{\binom{40}{2}}{2}\]is divided by $1000$.


 
 \textbf{Answer: }004



\textbf{Solution 1 (Telescoping)}

We first write the expression as a summation.
 \begin{align*} \sum_{i=3}^{40} \binom{\binom{i}{2}}{2}  & = \sum_{i=3}^{40} \binom{\frac{i \left( i - 1 \right)}{2}}{2} \\ & = \sum_{i=3}^{40} \frac{\frac{i \left( i - 1 \right)}{2} \left( \frac{i \left( i - 1 \right)}{2}- 1 \right)}{2} \\ & = \frac{1}{8} \sum_{i=3}^{40}  i \left( i - 1 \right) \left( i \left( i - 1 \right) - 2 \right)  \\ & = \frac{1}{8} \sum_{i=3}^{40} i(i - 1)(i^2-i-2) \\ & = \frac{1}{8} \sum_{i=3}^{40} i(i-1)(i+1)(i-2) \\ & = \frac{1}{8}\sum_{i=3}^{40} (i-2)(i-1)i(i+1) \\ & = \frac{1}{40}\sum_{i=3}^{40}[(i-2)(i-1)i(i+1)(i+2)-(i-3)(i-2)(i-1)i(i+1)]* \\ & = \frac{38\cdot39\cdot40\cdot41\cdot42-0}{40}\\ & = 38 \cdot 39 \cdot 41 \cdot 42 \\ & = \left( 40 - 2 \right) \left( 40 - 1 \right) \left( 40 + 1 \right) \left( 40 + 2 \right) \\ & = \left( 40^2 - 2^2 \right) \left( 40^2 - 1^2 \right) \\ & = \left( 40^2 - 4 \right) \left( 40^2 - 1 \right) \\ & = 40^4 - 40^2 \cdot 5 + 4 \\ & \equiv \boxed{004}\pmod{1000}\ \end{align*}
$*(i-2)(i-1)i(i+1)=\frac{1}{5}[(i-2)(i-1)i(i+1)(i+2)-(i-3)(i-2)(i-1)i(i+1)]$ is how we force the expression to telescope.~qyang

 





\textbf{Solution 2 (Hockey Stick)}

Doing simple algebra calculation will give the following equation:
 \begin{align*} \binom{\binom{n}{2}}{2} &= \frac{\frac{n(n-1)}{2} \cdot (\frac{n(n-1)}{2}-1)}{2}\\  &= \frac{n(n-1)(n^2-n-2)}{8}\\  &= \frac{(n+1)n(n-1)(n-2)}{8}\\  &= \frac{(n+1)!}{8\cdot (n-3)!} = 3 \cdot \frac{(n+1)!}{4!\cdot (n-3)!}\\  &= 3 \binom{n+1}{4} \end{align*}


 Next, by using \href{https://artofproblemsolving.com/wiki/index.php?title=Hockey-Stick\_Identity}{Hockey-Stick Identity}, we have:
 \[3 \cdot \sum_{i=3}^{40} \binom{i+1}{4} = 3 \binom{42}{5} = 42 \cdot 41 \cdot 39 \cdot 38\]
 \[=(40^2-2^2)(40^2-1^2) \equiv \boxed{004} ~(\text{mod}~ 1000)\]

 





\textbf{Solution 3}

Since $40$ seems like a completely arbitrary number, we can use Engineer's Induction by listing out the first few sums. These are, in the order of how many terms there are starting from $1$ term: $3$, $18$, $63$, $168$, $378$, and $756$. Notice that these are just $3 \cdot \dbinom50$, $3 \cdot \dbinom61$, $3 \cdot \dbinom72$, $3 \cdot \dbinom83$, $3 \cdot \dbinom94$, $3 \cdot \dbinom{10}5$. It's clear that this pattern continues up to $38$ terms, noticing that the "indexing" starts with $\dbinom32$ instead of $\dbinom12$. Thus, the value of the sum is $3 \cdot \dbinom{42}{37}=2552004 \equiv \boxed{\textbf{004}} \pmod{1000}$.

 





\textbf{Solution 4}

As in solution 1, obtain $\sum_{i=3}^{40} \binom{\binom{i}{2}}{2} = \frac{1}{8} \sum_{i=3}^{40} i^4-2i^3-i^2+2i.$ Write this as 
 \[\frac{1}{8}\left(\sum_{i=3}^{40} i^4 - 2\sum_{i=3}^{40}i^3 - \sum_{i=3}^{40}i^2 + 2\sum_{i=3}^{40}i\right).\]

 We can safely write this expression as $\frac{1}{8}\left(\sum_{i=1}^{40} i^4 - 2\sum_{i=1}^{40}i^3 - \sum_{i=1}^{40}i^2 + 2\sum_{i=1}^{40}i\right)$, since plugging $i=1$ and $i=2$ into $i^4-2i^3-i^2+2i$ both equal $0,$ meaning they won't contribute to the sum.

 Use the sum of powers formulae. We obtain 
 \[\frac{1}{8}\left(\frac{i(i+1)(2i+1)(3i^2+3i-1)}{30} - \frac{i^2(i+1)^2}{2} - \frac{i(i+1)(2i+1)}{6} + i(i+1)\right) \text{ where i = 40.}\]

 We can factor the following expression as $\frac{1}{8}\left(\frac{i(i+1)(2i+1)(3i^2+3i-6)}{30} - \frac{i(i+1)}{2} (i(i+1)-2)\right),$ and simplifying, we have 
 \[\sum_{i=3}^{40} \binom{\binom{i}{2}}{2} = \frac{i(i+1)(2i+1)(i^2+i-2)}{80}-\frac{i^2(i+1)^2-2i(i+1)}{16} \text{ where i = 40.}\]

 Substituting $i=40$ and simplifying gets $41\cdot 81\cdot 819 - 5\cdot 41\cdot 819,$ so we would like to find $819\cdot 76\cdot 41 \pmod{1000}.$ To do this, get $819\cdot 76\equiv 244 \pmod{1000}.$ Next, $244\cdot 41 \equiv \boxed{004} \pmod{1000}.$

 -sirswagger21

 


 





\textbf{Solution 5}

To solve this problem, we need to use the following result:

 
 \[ \sum_{i=n}^m \binom{i}{k} = \binom{m+1}{k+1} - \binom{n}{k+1} . \]

 Now, we use this result to solve this problem.

 We have
 \begin{align*} \sum_{i=3}^{40} \binom{\binom{i}{2}}{2}  & = \sum_{i=3}^{40} \binom{\frac{i \left( i - 1 \right)}{2}}{2} \\ & = \sum_{i=3}^{40} \frac{\frac{i \left( i - 1 \right)}{2} \left( \frac{i \left( i - 1 \right)}{2}- 1 \right)}{2} \\ & = \frac{1}{8} \sum_{i=3}^{40}  i \left( i - 1 \right) \left( i \left( i - 1 \right) - 2 \right)  \\ & = \frac{1}{8} \sum_{i=3}^{40}  i \left( i - 1 \right)  \left( \left( i - 2 \right) \left( i - 3 \right) + 4 \left( i - 2 \right) \right)  \\ & = 3 \left( \sum_{i=3}^{40} \binom{i}{4} + \sum_{i=3}^{40} \binom{i}{3} \right) \\ & = 3 \left( \binom{41}{5} - \binom{3}{5} + \binom{41}{4} - \binom{3}{4} \right) \\ & = 3 \left( \binom{41}{5} + \binom{41}{4} \right) \\ & = 3 \cdot \frac{41 \cdot 40 \cdot 39 \cdot 38}{5!} \left( 37 + 5 \right) \\ & = 3 \cdot 41 \cdot 13 \cdot 38 \cdot 42 \\ & = 38 \cdot 39 \cdot 41 \cdot 42 \\ & = \left( 40 - 2 \right) \left( 40 - 1 \right) \left( 40 + 1 \right) \left( 40 + 2 \right) \\ & = \left( 40^2 - 2^2 \right) \left( 40^2 - 1^2 \right) \\ & = \left( 40^2 - 4 \right) \left( 40^2 - 1 \right) \\ & = 40^4 - 40^2 \cdot 5 + 4  . \end{align*}


 Therefore, modulo 1000, $\sum_{i=3}^{40} \binom{\binom{i}{2}}{2}  \equiv \boxed{\textbf{(004) }}$.

 





\textbf{Solution 6 (Combinatorial Method)}

We examine the expression $\binom{\binom{n}{2}}{2}$. Imagine we have a set $S$ of $n$ integers. Then the expression can be translated to the number of pairs of $2$ element subsets of $S$.

 To count this, note that each pair of $2$ element subsets can either share $1$ value or $0$ values. In the former case, pick three integers $a$, $b$, and $c$. There are $\binom{n}{3}$ ways to select these integers and $3$ ways to pick which one of the three is the shared integer. This gives $3\cdot \binom{n}{3}$.

 In the latter case, we pick $4$ integers $a$, $b$, $c$, and $d$ in a total of $\binom{n}{4}$ ways. There are $\frac{1}{2}\binom{4}{2} = 3$ ways to split this up into $2$ sets of $2$ integers — $\binom{4}{2}$ ways to pick which $2$ integers are together and dividing by $2$ to prevent overcounting. This gives $3\cdot \binom{n}{4}$.

 So we have 
 \[\binom{\binom{n}{2}}{2} = 3\cdot \binom{n}{3} + 3\cdot \binom{n}{4}\]We use the Hockey Stick Identity to evaluate this sum:
 \begin{align*} \sum_{n=3}^{40} \binom{\binom{n}{2}}{2} &= \sum_{n=3}^{40} \left( 3 \binom{n}{3} + 3 \binom{n}{4} \right) \\ &= 3 \left( \sum_{n=3}^{40}  \binom{n}{3} \right) + 3\left( \sum_{n=4}^{40}  \binom{n}{4} \right) \\ &= 3\left( \binom{41}{4} + \binom{41}{5} \right) \end{align*}
Evaluating while accounting for mod $1000$ gives the final answer to be $\boxed{004}$.

 
Note that $\binom{n}{3}+\binom{n}{4}=\binom{n+1}{4}$ (Pascal's identity), so this result is consistent with other solutions. ~TThB0501

 


 


 





\textbf{Solution 6 (elegant)}

think about an n-gon. nc2= amount of ways to make an edge. (nc2)c2=amount of ways to pick two random edges out of these edges. however, we can reconstruct this. since no edge will be picked twice, there will either be 3 points or 4 points used in the edges. if there are 3 points, we can use nc3 to outline a triangle, and multiply 3 for the 3 options of sides we could remove to make two edges intersecting at one point. if there are 4 points, we note that given any 4 points, just using those configurations, there are only 3 possible ways to make edges with all points. Therefore, we can say that (nc2)c2=3(nc3)+3(nc4).

 then go as above

 




	\item (\textit{AIME II 2022}) Let $ABCD$ be a convex quadrilateral with $AB=2, AD=7,$ and $CD=3$ such that the bisectors of acute angles $\angle{DAB}$ and $\angle{ADC}$ intersect at the midpoint of $\overline{BC}.$ Find the square of the area of $ABCD.$


 
 \textbf{Answer: }180



\textbf{Solution 1}


                    \begin{figure}[H]
                        \centering
                        \includegraphics[width=0.5\textwidth]{Solution/AIME_2022_II_11_Solution_Figure_1.png}
                    \end{figure}
                    

 According to the problem, we have $AB=AB'=2$, $DC=DC'=3$, $MB=MB'$, $MC=MC'$, and $B'C'=7-2-3=2$

 Because $M$ is the midpoint of $BC$, we have $BM=MC$, so: 
 \[MB=MB'=MC'=MC.\]

 Then, we can see that $\bigtriangleup{MB'C'}$ is an isosceles triangle with $MB'=MC'$

 Therefore, we could start our angle chasing: $\angle{MB'C'}=\angle{MC'B'}=180^\circ-\angle{MC'D}=180^\circ-\angle{MCD}$.

 This is when we found that points $M$, $C$, $D$, and $B'$ are on a circle. Thus, $\angle{BMB'}=\angle{CDC'} \Rightarrow \angle{B'MA}=\angle{C'DM}$. This is the time we found that $\bigtriangleup{AB'M} \sim \bigtriangleup{MC'D}$.

 Thus, $\frac{AB'}{B'M}=\frac{MC'}{C'D} \Longrightarrow (B'M)^2=AB' \cdot C'D = 6$

 Point $H$ is the midpoint of $B'C'$, and $MH \perp AD$. $B'H=HC'=1 \Longrightarrow MH=\sqrt{B'M^2-B'H^2}=\sqrt{6-1}=\sqrt{5}$.

 The area of this quadrilateral is the sum of areas of triangles: 
 \[S_{\bigtriangleup{ABM}}+S_{\bigtriangleup{AB'M}}+S_{\bigtriangleup{CDM}}+S_{\bigtriangleup{C'DM}}+S_{\bigtriangleup{B'C'M}}\]
 \[=S_{\bigtriangleup{AB'M}}\cdot 2 + S_{\bigtriangleup{B'C'M}} + S_{\bigtriangleup{C'DM}}\cdot 2\]
 \[=2 \cdot \frac{1}{2} \cdot AB' \cdot MH + \frac{1}{2} \cdot B'C' \cdot MH + 2 \cdot \frac{1}{2} \cdot C'D \cdot MH\]
 \[=2\sqrt{5}+\sqrt{5}+3\sqrt{5}=6\sqrt{5}\]

 Finally, the square of the area is $(6\sqrt{5})^2=\boxed{180}$

 





\textbf{Solution 2}

Denote by $M$ the midpoint of segment $BC$.Let points $P$ and $Q$ be on segment $AD$, such that $AP = AB$ and $DQ = DC$.

 Denote $\angle DAM = \alpha$, $\angle BAD = \beta$, $\angle BMA = \theta$, $\angle CMD = \phi$.

 Denote $BM = x$. Because $M$ is the midpoint of $BC$, $CM = x$.

 Because $AM$ is the angle bisector of $\angle BAD$ and $AB = AP$, $\triangle BAM \cong \triangle PAM$.Hence, $MP = MB$ and $\angle AMP = \theta$.Hence, $\angle MPD = \angle MAP + \angle PMA = \alpha + \theta$.

 Because $DM$ is the angle bisector of $\angle CDA$ and $DC = DQ$, $\triangle CDM \cong \triangle QDM$.Hence, $MQ = MC$ and $\angle DMQ = \phi$.Hence, $\angle MQA = \angle MDQ + \angle QMD = \beta + \phi$.

 Because $M$ is the midpoint of segment $BC$, $MB = MC$.Because $MP = MB$ and $MQ = MC$, $MP = MQ$.

 Thus, $\angle MPD = \angle MQA$.

 Thus,
 \[ \alpha + \theta = \beta + \phi . \hspace{1cm} (1) \]

 In $\triangle AMD$, $\angle AMD = 180^\circ - \angle MAD - \angle MDA = 180^\circ - \alpha - \beta$.In addition, $\angle AMD = 180^\circ - \angle BMA - \angle CMD = 180^\circ - \theta - \phi$.Thus,
 \[ \alpha + \beta = \theta + \phi . \hspace{1cm} (2) \]

 Taking $(1) + (2)$, we get $\alpha = \phi$.Taking $(1) - (2)$, we get $\beta = \theta$.

 Therefore, $\triangle ADM \sim \triangle AMB \sim \triangle MDC$.

 Hence, $\frac{AD}{AM} = \frac{AM}{AB}$ and $\frac{AD}{DM} = \frac{DM}{CD}$.Thus, $AM = \sqrt{AD \cdot AD} = \sqrt{14}$ and $DM = \sqrt{AD \cdot CD} = \sqrt{21}$.

 In $\triangle ADM$, by applying the law of cosines, $\cos \angle AMD  = \frac{AM^2 + DM^2 - AD^2}{2 AM \cdot DM} = - \frac{1}{\sqrt{6}}$.Hence, $\sin \angle AMD = \sqrt{1 - \cos^2 \angle AMD} = \frac{\sqrt{5}}{\sqrt{6}}$.Hence, ${\textrm Area} \ \triangle ADM = \frac{1}{2} AM \cdot DM \dot \sin \angle AMD = \frac{7 \sqrt{5}}{2}$.

 Therefore,
 \begin{align*} {\textrm Area} \ ABCD & = {\textrm Area} \ \triangle AMD + {\textrm Area} \ \triangle ABM + {\textrm Area} \ \triangle MCD \\ & = {\textrm Area} \ \triangle AMD \left( 1 + \left( \frac{AM}{AD} \right)^2 + \left( \frac{MD}{AD} \right)^2 \right) \\ & = 6 \sqrt{5} . \end{align*}


 Therefore, the square of ${\textrm Area} \ ABCD$ is $\left( 6 \sqrt{5} \right)^2 = \boxed{\textbf{(180) }}$.

 





\textbf{Solution 3 (Visual)}


                    \begin{figure}[H]
                        \centering
                        \includegraphics[width=0.5\textwidth]{Solution/AIME_2022_II_11_Solution_Figure_2.png}
                    \end{figure}
                    \textbf{\textit{Claim}}

 In the triangle $ABC, AB = 2AC, M$ is the midpoint of $AB. D$ is the point of intersection of the circumcircle and the bisector of angle $A.$ Then $DM = BD.$

 \textbf{\textit{ Proof}}

 Let $A = 2\alpha.$ Then $\angle DBC = \angle DCB = \alpha.$

 Let $E$ be the intersection point of the perpendicular dropped from $D$ to $AB$ with the circle.

 Then the sum of arcs  $\overset{\Large\frown} {BE} + \overset{\Large\frown}{AC} + \overset{\Large\frown}{CD} = 180^\circ.$
 \[\overset{\Large\frown} {BE} = 180^\circ – 2\alpha – \overset{\Large\frown}{AC}.\]

 Let $E'$ be the point of intersection of the line $CM$ with the circle.$CM$ is perpendicular to $AD, \angle AMC = 90^\circ – \alpha,$ the sum of arcs  $\overset{\Large\frown}{A}C +  \overset{\Large\frown}{BE'} = 180^\circ – 2\alpha \implies E'$ coincides with $E.$

 The inscribed angles $\angle DEM = \angle DEB, M$ is symmetric to $B$ with respect to $DE, DM = DB.$

 \textbf{\textit{ Solution}}

 
                    \begin{figure}[H]
                        \centering
                        \includegraphics[width=0.5\textwidth]{Solution/AIME_2022_II_11_Solution_Figure_3.png}
                    \end{figure}
                    Let $AB' = AB, DC' = DC, B'$ and $C'$ on $AD.$

 Then $AB' = 2, DC' = 3, B'C' = 2 = AB'.$

 Quadrilateral  $ABMC'$ is cyclic.Let $\angle A = 2\alpha.$ Then $\angle MBC' = \angle MC'B = \alpha.$

 Circle $BB'C'C$ centered at $M, BC$ is its diameter, $\angle BC'C = 90^\circ.$$\angle DMC' = \angle MC'B,$ since they both complete $\angle MC'C$ to $90^\circ.$

 $\angle MB'A = \angle MC'D,$ since they are the exterior angles of an isosceles $\triangle MB'C'.$$\triangle AMB' \sim \triangle MDC'$ by two angles.$\frac {AB'}{MC'} = \frac {MB'}{DC'}, MC' =\sqrt{AB' \cdot C'D} = \sqrt{6}.$

 The height dropped from $M$ to $AD$ is $\sqrt{MB'^2 - (\frac{B'C'}{2})^2} =\sqrt{6 - 1} = \sqrt{5}.$

 The areas of triangles $\triangle AMB'$ and $\triangle MC'B'$ are equal to $\sqrt{5},$ area of $\triangle MC'D$ is $\frac{3}{2} \sqrt{5}.$

 
 \[\triangle AMB' = \triangle AMB, \triangle MC'D = \triangle MCD \implies\] The area of  $ABCD$ is $(1 + 2 + 3) \sqrt{5} = 6\sqrt{5} \implies 6^2 \cdot 5 = \boxed{180}.$

 \textbf{vladimir.shelomovskii@gmail.com, vvsss}

 





\textbf{Solution 4 (THINK OUTSIDE THE BOX)}

Extend $AB$ and $CD$ so they intersect at a point $X$. Then note that $M$ is the incenter of $\triangle{XAD}$, implying that $M$ is on the angle bisector of $X$. Now because $XM$ is both an angle bisector and a median of $\triangle{XBC}$, $\triangle{XBC}$ is isosceles. Then we can start angle chasing:

 Let $\angle{BAM}=a, \angle{CDM}=b,$ and $\angle{XBC}=c$. Then $\angle{AMD}=\pi-(a+b), \angle{ABM}=\pi-c, \angle{DCM}=\pi-c$, implying that $\angle{BMA}+\angle{CMD}=a+b$, implying that $2c-(a+b)=(a+b)$, or that $c=a+b$. Substituting this into the rest of the diagram, we find that $\triangle{BMA} \sim \triangle{CDM} \sim \triangle{MDA}$. 

 Then $\frac{AB}{BM}=\frac{MC}{CD}$, or $BM=CM=\sqrt{6}$. Moreover, $\frac{AB}{AM}=\frac{AM}{AD}$, or $AM=\sqrt{14}$. Similarly, $\frac{CD}{MD}=\frac{MD}{AD}$, or $DM=\sqrt{21}$. Then using Law of Cosines on $\triangle{AMD}$, to get that $cos\angle{AMD}=-\frac{\sqrt{6}}{6}$, or $sin\angle{AMD}=\frac{\sqrt{30}}{6}$. 

 We finish by using the formula $K=\frac{1}{2}absinC$, as follows:

 $[ABCD]=[ABM]+[CDM]+[ADM]=\frac{\frac{\sqrt{30}}{6}(2\sqrt{6}+3\sqrt{6}+7\sqrt{6})}{2}=6\sqrt{5}$.

 $(6\sqrt{5})^2=\boxed{180}$.

 -dragoon

 





\textbf{Solution 5 (incircle)}


                    \begin{figure}[H]
                        \centering
                        \includegraphics[width=0.5\textwidth]{Solution/AIME_2022_II_11_Solution_Figure_4.png}
                    \end{figure}
                    As shown in paragraph one of solution 4, extending $AB$ and $CD$ to $X$, we realize that $\triangle{XBC}$ is isosceles, thus $XM \perp BC$. Let $XB = XC = x$. And, midpoint $M$ is the incenter of $\triangle{XAD}$. Construct perpendiculars $ME, MF, MG$ to sides $AD, AX, DX$ respectively (constructing the radii of the incircle). Let $EM = FM = GM = r$. The semiperimeter $s = \frac{2x + 2 + 3 + 7}{2} = x+6$. Since $FX$ is the tangent off the incircle, $FX = s - AD = x-1$. So, $BF = BX - FX = 1$. Because $\triangle{BFM} \sim \triangle{MFX}$,

 
 \[\frac{BF}{FM} = \frac{FM}{FX} \implies {FM}^2 = BF \cdot FX = x - 1 \implies r^2 = x - 1\].

 By Heron's formula and the inradius area formula,

 
 \[(x+6)r = \sqrt{(x+6)\cdot 4 \cdot 3 \cdot (x-1)} \implies (x+6)r^2 = 12(x-1) \implies (x+6)(x-1) = 12(x-1) \implies x=6\]

 Then,  $r^2 = x - 1 = 5 \implies r = \sqrt{5}$. Finally,

 
 \[[ABCD] = [ABM] + [CDM] + [AMD] = \frac{AB \cdot FM}{2} + \frac{CD \cdot GM}{2} + \frac{AD \cdot EM}{2} = \frac{2r}{2} + \frac{3r}{2} + \frac{7r}{2} = 6r = 6\sqrt{5}\]

 Thus, our answer is $(6\sqrt{5})^{2} = \boxed{180}$

 \textbf{Remark} $B$ and $C$ are the tangent points of the X-mixtilinear incircle (of $\triangle{XAD}$). This may be useful, but I haven't looked into it.

 Note:You could also derive $r^2 = x - 1$ by using Geometric Mean Theorem as triangle $BMX$ is right.~anishsaichalam

 





\textbf{Solution 6 (bash)}

Let the midpoint of $BC$ be $M$. Angle-chase and observe that $\Delta AMD~\Delta ABM~\Delta MCD$. Let $BM=CM=a$ and $AM=x$ and $DM=y$. As a result of this similarity, we write 

 
 \[\dfrac2a=\dfrac a3,\]

 which gives $a=\sqrt 6$. Similarly, we write 

 
 \[\dfrac2x=\dfrac x7\]

 and 

 
 \[\dfrac3y=\dfrac y7\]

 to get $x=\sqrt{14}$ and $y=\sqrt{21}$. 

 We now have all required side lengths; we can find the area of $\Delta AMD$ with Heron's formula. Doing so yields $\dfrac72\sqrt5$. We could also bash out the areas of the other two triangles since we know all their side lengths (this is what I did :sob:), but a more intelligent method is to recall the triangles' similarity. The ratio of similarlity between $\Delta AMD$ and $\Delta ABM$ is $\dfrac{\sqrt{14}}7=\sqrt{\dfrac27}$, and between $\Delta AMD$ and $\Delta MCD$ is $\dfrac{\sqrt{21}}7=\sqrt{\dfrac37}$. Thus, the area ratios are $\dfrac27$ and $\dfrac37$, respectively, so adding together we have $\dfrac27+\dfrac37=\dfrac57$. Multiplying this by our $\dfrac72\sqrt5$, we have $\dfrac52\sqrt5$ as their total area. Adding this to our original area, we have $\dfrac52\sqrt5+\dfrac72\sqrt5=\sqrt5\left(\dfrac52+\dfrac72\right)=\sqrt5\left(\dfrac{12}2\right)=6\sqrt5$. 

 The square of this is $\boxed{180}$. 

 





\textbf{Solution 7 (similar to 4)}

As in solution 4, let $X=AB\cap CD$, so $M$ is the incenter of $ADX$ and $XB=XC$. Let $XB=XC=x$. Then the normalized barycentric coordinates of $B$, $C$, and $M$ with respect to $ADX$ are $\left[\frac{x}{x+2}:0:\frac{2}{x+2}\right]$, $\left[0:\frac{x}{x+3}:\frac{3}{x+3}\right]$, and $\left[\frac{x+3}{2x+12}:\frac{x+2}{2x+12}:\frac{7}{2x+12}\right]$. So we have $\frac{1}{2}\frac{x}{x+2}=\frac{x+3}{2x+12}$ giving $x=6$. The sidelengths of $ADX$ are thus $AD=7$, $AX=8$, and $DX=9$ giving $[ADX]=12\sqrt 5$. Also, we have $[BCX]=\frac{6}{8}\cdot\frac{6}{9}[ADX]=6\sqrt 5$ so that $[ABCD]=[ADX]-[BCX]=6\sqrt 5$. The area squared is thus $\boxed{180}$.~~ golue3120

 




	\item (\textit{AIME II 2022}) Let $a, b, x,$ and $y$ be real numbers with $a>4$ and $b>1$ such that
 \[\frac{x^2}{a^2}+\frac{y^2}{a^2-16}=\frac{(x-20)^2}{b^2-1}+\frac{(y-11)^2}{b^2}=1.\]Find the least possible value of $a+b.$


 
 \textbf{Answer: }023



\textbf{Solution}

Denote $P = \left( x , y \right)$.

 Because $\frac{x^2}{a^2}+\frac{y^2}{a^2-16} = 1$, $P$  is on an ellipse whose center is $\left( 0 , 0 \right)$ and foci are $\left( - 4 , 0 \right)$ and $\left( 4 , 0 \right)$.

 Hence, the sum of distance from $P$ to $\left( - 4 , 0 \right)$ and $\left( 4 , 0 \right)$ is equal to twice the major axis of this ellipse, $2a$.

 Because $\frac{(x-20)^2}{b^2-1}+\frac{(y-11)^2}{b^2} = 1$, $P$  is on an ellipse whose center is $\left( 20 , 11 \right)$ and foci are $\left( 20 , 10 \right)$ and $\left( 20 , 12 \right)$.

 Hence, the sum of distance from $P$ to $\left( 20 , 10 \right)$ and $\left( 20 , 12 \right)$ is equal to twice the major axis of this ellipse, $2b$.

 Therefore, $2a + 2b$ is the sum of the distance from $P$ to four foci of these two ellipses.

 To make this minimized, $P$ is the intersection point of the line that passes through $\left( - 4 , 0 \right)$ and $\left( 20 , 10 \right)$, and the line that passes through $\left( 4 , 0 \right)$ and $\left( 20 , 12 \right)$.

 The distance between $\left( - 4 , 0 \right)$ and $\left( 20 , 10 \right)$ is $\sqrt{\left( 20 + 4 \right)^2 + \left( 10 - 0 \right)^2} = 26$.

 The distance between $\left( 4 , 0 \right)$ and $\left( 20 , 12 \right)$ is $\sqrt{\left( 20 - 4 \right)^2 + \left( 12 - 0 \right)^2} = 20$.

 Hence, $2 a + 2 b = 26 + 20 = 46$.

 Therefore, $a + b = \boxed{\textbf{(023) }}.$

 


 




	\item (\textit{AIME II 2022}) There is a polynomial $P(x)$ with integer coefficients such that
 \[P(x)=\frac{(x^{2310}-1)^6}{(x^{105}-1)(x^{70}-1)(x^{42}-1)(x^{30}-1)}\]holds for every $0<x<1.$ Find the coefficient of $x^{2022}$ in $P(x)$.


 
 \textbf{Answer: }220



\textbf{Solution 1}

Because $0 < x < 1$, we have
 \begin{align*} P \left( x \right) & = \sum_{a=0}^6  \sum_{b=0}^\infty \sum_{c=0}^\infty \sum_{d=0}^\infty \sum_{e=0}^\infty \binom{6}{a} x^{2310a} \left( - 1 \right)^{6-a} x^{105b} x^{70c} x^{42d} x^{30e} \\ & = \sum_{a=0}^6 \sum_{b=0}^\infty \sum_{c=0}^\infty \sum_{d=0}^\infty \sum_{e=0}^\infty \left( - 1 \right)^{6-a} x^{2310 a + 105 b + 70 c + 42 d + 30 e} . \end{align*}


 Denote by $c_{2022}$ the coefficient of $P \left( x \right)$.Thus,
 \begin{align*} c_{2022} & = \sum_{a=0}^6 \sum_{b=0}^\infty \sum_{c=0}^\infty \sum_{d=0}^\infty \sum_{e=0}^\infty \left( - 1 \right)^{6-a} \Bbb I \left\{ 2310 a + 105 b + 70 c + 42 d + 30 e = 2022 \right\} \\ & =  \sum_{b=0}^\infty \sum_{c=0}^\infty \sum_{d=0}^\infty \sum_{e=0}^\infty \left( - 1 \right)^{6-0} \Bbb I \left\{ 2310 \cdot 0 + 105 b + 70 c + 42 d + 30 e = 2022 \right\} \\ & = \sum_{b=0}^\infty \sum_{c=0}^\infty \sum_{d=0}^\infty \sum_{e=0}^\infty \Bbb I \left\{ 105 b + 70 c + 42 d + 30 e = 2022 \right\} . \end{align*}


 Now, we need to find the number of nonnegative integer tuples $\left( b , c , d , e \right)$ that satisfy 
 \[ 105 b + 70 c + 42 d + 30 e = 2022 . \hspace{1cm} (1) \]

 Modulo 2 on Equation (1), we have $b \equiv 0 \pmod{2}$.Hence, we can write $b = 2 b'$. Plugging this into (1), the problem reduces to finding the number of nonnegative integer tuples $\left( b' , c , d , e \right)$ that satisfy
 \[ 105 b' + 35 c + 21 d + 15 e = 1011 . \hspace{1cm} (2) \]

 Modulo 3 on Equation (2), we have $2 c \equiv 0 \pmod{3}$.Hence, we can write $c = 3 c'$. Plugging this into (2), the problem reduces to finding the number ofnonnegative integer tuples $\left( b' , c' , d , e \right)$ that satisfy
 \[ 35 b' + 35 c' + 7 d + 5 e = 337 . \hspace{1cm} (3) \]

 Modulo 5 on Equation (3), we have $2 d \equiv 2 \pmod{5}$.Hence, we can write $d = 5 d' + 1$. Plugging this into (3), the problem reduces to finding the number ofnonnegative integer tuples $\left( b' , c' , d' , e \right)$ that satisfy
 \[ 7 b' + 7 c' + 7 d' + e = 66 . \hspace{1cm} (4) \]

 Modulo 7 on Equation (4), we have $e \equiv 3 \pmod{7}$.Hence, we can write $e = 7 e' + 3$. Plugging this into (4), the problem reduces to finding the number ofnonnegative integer tuples $\left( b' , c' , d' , e' \right)$ that satisfy
 \[ b' + c' + d' + e' = 9 . \hspace{1cm} (5) \]

 The number of nonnegative integer solutions to Equation (5) is $\binom{9 + 4 - 1}{4 - 1} = \binom{12}{3} =  \boxed{\textbf{(220) }}$.

 





\textbf{Solution 2}

Note that $2022 = 210\cdot 9 +132$. Since the only way to express $132$ in terms of $105$, $70$, $42$, or $30$ is $132 = 30+30+30+42$, we are essentially just counting the number of ways to express $210*9$ in terms of these numbers. Since $210 = 2*105=3*70=5*42=7*30$, it can only be expressed as a sum in terms of only one of the numbers ($105$, $70$, $42$, or $30$). Thus, the answer is (by sticks and stones/stars and bars) 
 \[\binom{12}{3} = \boxed{\textbf{(220)}}\]

 





\textbf{Solution 3}

We know that $\frac{a^n-b^n}{a-b}=\sum_{i=0}^{n-1} a^{n-1-i}b^i$. Applying this, we see that 
 \[P(x)=(1+x^{105}+x^{210}+...)(1+x^{70}+x^{140}+...)(1+x^{42}+x^{84}+...)(1+x^{30}+x^{60}+...)(x^{4620}-2x^{2310}+1)\] The last factor does not contribute to the $x^{2022}$ term, so we can ignore it. Thus we only have left to solve the equation $105b+70c+42d+30e=2022$, and we can proceed from here with Solution 1.

 





\textbf{Solution 4 (Generating Function Bash)}

Essentially we want the coefficient of $x^{2022}$ in the expansion

 
 \[\frac{1}{(1 - x^{105})(1 - x^{70})(1 - x^{42})(1 - x^{30})}\]

 As $(x^{2310} -1)^{6}$ does not contribute to the expansion, we omit it.

 Notice $\text{lcm}(105,70,42,30) = 210$, so we can rewrite the generating function as

 

 \[\frac{(1 + x^{105})(1 + x^{70} + x^{140})(1 + x^{42} + ... + x^{168})(1 + x^{30} + ... + x^{180})}{(1-x^{210})^{4}}\]

 Notice to obtain a $x^{2022}$ value, the $x^{42}$ must be used, so we can reduce the problem to finding the coefficient of $x^{1980}$ in 

 
 \[\frac{(1 + x^{105})(1 + x^{70} + x^{140})(1 + x^{30} + ... + x^{180})}{(1-x^{210})^{4}}\]

 If $(1 - x^{210})^{-4}$ is expanded, it will only generate multiples of 210. To compensate this, notice $1980 \equiv 90 \pmod{210}$, which means we need terms with a power of a 90 to acheive what we want (300 and larger values can be shown to be impossible upon inspection). This implies that the first two brackets do not contribute and we are left with

 
 \[\frac{x^{90}}{(1-x^{210})^{4}}\]

 This reduces to finding the coefficient of $x^{1890}$ for the expansion $(1 - x^{210})^{-4}$, which is $\binom{12}{3} = \fbox{220}$

 





\textbf{Solution 5 (Similar to solution 1)}

Expand the numerator and divide each term in the denominator separately to get 

 \begin{verbatim}
	$(1+x^{30}+x^{60}+\dots)\cdot         (1+x^{42}+x^{84}+\dots)\cdot(1+x^{70}+x^{140}+\dots)$     $\cdot (1+x^{105}+x^{210}+\dots) (x^{2310}-1)^{2}$
\end{verbatim}Since the factor of $(x^{2310}-1)^{2}$ is too big so we choose the -1 from both the 2 factors to get a 1.So now we need to find the number of quadruples ofnon negative integers $(a,b,c,d)$ such that$\textbf{105a + 70b + 42c + 30d = 2022}$.By considering modulo $2$, $3$, $5$, $7$ we can reduce the above expression to something that can be solved using stars-and-bars. 

 \begin{verbatim}
	$a \equiv 0 \pmod 2$ $b \equiv 0 \pmod 3$ $c \equiv 1 \pmod 5$ $d \equiv 3 \pmod7$.
\end{verbatim}Now we put in $a=2x$, $b=3y$, $c=5z+1$ and $d=7w+3$ to get

 \begin{verbatim}
$2022 = 210(x+y+z+w) + 132$   $\Rightarrow \boxed{x+y+z+w=9},\; x,y,z,w\geq0$
\end{verbatim}By stars-and-bars the answer is $\binom{9+4-1}{4-1} = \boxed{\textbf{220}}$.

 




	\item (\textit{AIME II 2022}) For positive integers $a$, $b$, and $c$ with $a < b < c$, consider collections of postage stamps in denominations $a$, $b$, and $c$ cents that contain at least one stamp of each denomination. If there exists such a collection that contains sub-collections worth every whole number of cents up to $1000$ cents, let $f(a, b, c)$ be the minimum number of stamps in such a collection. Find the sum of the three least values of $c$ such that $f(a, b, c) = 97$ for some choice of $a$ and $b$.


 
 \textbf{Answer: }188



\textbf{Solution}

Notice that we must have $a = 1$; otherwise $1$ cent stamp cannot be represented. At least $b-1$ numbers of $1$ cent stamps are needed to represent the values less than $b$. Using at most $c-1$ stamps of value $1$ and $b$, it can have all the values from $1$ to $c-1$ cents. Plus $\lfloor \frac{999}{c} \rfloor$ stamps of value $c$, every value up to $1000$ can be represented. 

 \textbf{Correction:} This should be $\lfloor \frac{1000}{c} \rfloor$. The current function breaks when $c \mid 1000$ and $b \mid c$. Take $c = 200$ and $b = 20$. Then, we have $\lfloor \frac{999}{200} \rfloor = 4$ stamps of value 200, $\lfloor \frac{199}{20} \rfloor = 9$ stamps of value b, and 19 stamps of value 1. The maximum such a collection can give is $200 \cdot 4 + 20 \cdot 9 +19 \cdot 1 = 999$, just shy of the needed 1000. As for the rest of solution, proceed similarly, except use $1000$ instead of $999$.

 Also, some explanation: $b-1$ one cent stamps cover all residues module $b$. Having $\lfloor \frac{c-1}{b} \rfloor$ stamps of value b covers all residue classes modulo $c$. Finally, we just need $\lfloor \frac{1000}{c} \rfloor$ to cover everything up to 1000.

 In addition, note that this function sometimes may not always minimize the number of stamps required. This is due to the fact that the stamps of value $b$ and of value $1$ have the capacity to cover values greater than or equal to $c$ (which occurs when $c-1$ has a remainder less than $b-1$ when divided by $b$). Thus, in certain cases, not all $\lfloor \frac{1000}{c} \rfloor$ stamps of value c may be necessary, because the stamps of value $b$ and 1 can replace one $c$.

 


 \textbf{Correction for the correction and the original solution:} Actually, $f(a, b, c) = b + \lceil \frac{c}{b} \rceil + \lceil \frac{1001-b \lceil \frac{c}{b} \rceil}{c} \rceil - 2$.

 This could be obtained by solving the minimum of $A+B+C$ for the system of inequalities $A \geq b-14$, $A+Bb \geq c$, and $A+Bb+Cc \geq 1000$, (Using the inequality adjustment method) where $A$, $B$, and $C$ represent the number of $a$, $b$, and $c$ postage stamps used. It is possible to get that $A=b-1$, $B=\lceil \frac{c}{b} \rceil-1$, $C=\lceil \frac{1001-b \lceil \frac{c}{b} \rceil}{c} \rceil$. $f(a,b,c)=A+B+C$ yields the above.

 Note that in $\text{Case } 2$, $c=88$, for the original solution $a=1$, $b=86$, $c=88$, $f(a,b,c)=96$ if the above definition is used.

 When $A=85$, $B=1$, and $C=10$, you can get a collection that satisfies the conditions of the problem. $A+B+C=96$ here.

 Note that using the correction does not change that the original $f(1,87,88)=97$.

 Under the correction of the correction:

 The actual solution could be achieved by using inequalities on the left and right of $f(a,b,c)$ to get

 $b+\frac{1001}{c}-3<f(a,b,c)<c+1+\frac{1001}{c}$.

 Similar steps follow after substituting $f(a,b,c)=97$.

 However, $f(1,7,11)=f(1,87,88)=f(1,87,89)=97$.

 Therefore the answer is still $11+88+89=\boxed{188}$.

 ————————————————————————————————————

 
Therefore using $\lfloor \frac{999}{c} \rfloor$ stamps of value $c$, $\lfloor \frac{c-1}{b} \rfloor$ stamps of value $b$, and $b-1$ stamps of value $1$, all values up to $1000$ can be represented in sub-collections, while minimizing the number of stamps.

 So, $f(a, b, c) = \lfloor \frac{999}{c} \rfloor + \lfloor \frac{c-1}{b} \rfloor + b-1$.

 $\lfloor \frac{999}{c} \rfloor + \lfloor \frac{c-1}{b} \rfloor + b-1 = 97$. We can get the answer by solving this equation.

 $c > \lfloor \frac{c-1}{b} \rfloor + b-1$

 $\frac{999}{c} + c > \lfloor \frac{999}{c} \rfloor + \lfloor \frac{c-1}{b} \rfloor + b-1 = 97$

 $c^2 - 97c + 999 > 0$, $c > 85.3$ or $c < 11.7$

 $\lfloor \frac{999}{c} \rfloor + \lfloor \frac{c-1}{b} \rfloor + b-1 > \frac{999}{c}$

 $97 > \frac{999}{c}$, $c>10.3$

 \begin{verbatim}
$\text{Case } 1:$ $10.3 < c < 11.7 \rightarrow c = 11 \rightarrow \lfloor \frac{999}{11} \rfloor + \lfloor \frac{10}{b} \rfloor + b-1 = 97$$\lfloor \frac{10}{b} \rfloor + b = 8$, $b=7$
\end{verbatim}\begin{verbatim}
$\text{Case } 2: c>85.3$
\end{verbatim}\begin{verbatim}
$c = 86 \rightarrow$ $\lfloor \frac{999}{86} \rfloor + \lfloor \frac{85}{b} \rfloor + b-1 = 97$$\lfloor \frac{85}{b} \rfloor + b = 87$, $b=87 > c \rightarrow \text{no solution}$
\end{verbatim}\begin{verbatim}
$c = 87 \rightarrow$ $\lfloor \frac{999}{87} \rfloor + \lfloor \frac{86}{b} \rfloor + b-1 = 97$$\lfloor \frac{86}{b} \rfloor + b = 87$, $b=86$ or $1$. $\text{We cannot have }b=1 \text{ since it doesn't satisfy } a<b \text{, and note that if } b=86 \text{ we can have 10 coins of value } c, 1 \text{ of } b, \text{ and } 85$ $\text{of } a \text{ for a total of 96 coins and still be able to make every value from 1 to 1000. Thus } c=87 \text{ yields no solution.}$
\end{verbatim}\begin{verbatim}
$c = 88 \rightarrow$ $\lfloor \frac{999}{88} \rfloor + \lfloor \frac{87}{b} \rfloor + b-1 = 97$$\lfloor \frac{87}{b} \rfloor + b = 87$, $b=86$
\end{verbatim}\begin{verbatim}
$c = 89 \rightarrow$ $\lfloor \frac{999}{89} \rfloor + \lfloor \frac{88}{b} \rfloor + b-1 = 97$$\lfloor \frac{88}{b} \rfloor + b = 87$, $b=86$
\end{verbatim}The $3$ least values of $c$ are $11$, $88$, $89$. $11 + 88+ 89 = \boxed{\textbf{188}}$

 


 




	\item (\textit{AIME II 2022}) Two externally tangent circles $\omega_1$ and $\omega_2$ have centers $O_1$ and $O_2$, respectively. A third circle $\Omega$ passing through $O_1$ and $O_2$ intersects $\omega_1$ at $B$ and $C$ and $\omega_2$ at $A$ and $D$, as shown. Suppose that $AB = 2$, $O_1O_2 = 15$, $CD = 16$, and $ABO_1CDO_2$ is a convex hexagon. Find the area of this hexagon.

                    \begin{figure}[H]
                        \centering
                        \includegraphics[width=0.5\textwidth]{Problem/AIME_2022_II_15_Problem_Figure_1.png}
                    \end{figure}
                    


 
 \textbf{Answer: }140



\textbf{Solution 1}

First observe that $AO_2 = O_2D$ and $BO_1 = O_1C$. Let points $A'$ and $B'$ be the reflections of $A$ and $B$, respectively, about the perpendicular bisector of $\overline{O_1O_2}$. Then quadrilaterals $ABO_1O_2$ and $B'A'O_2O_1$ are congruent, so hexagons $ABO_1CDO_2$ and $A'B'O_1CDO_2$ have the same area. Furthermore, triangles $DO_2A'$ and $B'O_1C$ are congruent, so $A'D = B'C$ and quadrilateral $A'B'CD$ is an isosceles trapezoid.
                    \begin{figure}[H]
                        \centering
                        \includegraphics[width=0.5\textwidth]{Solution/AIME_2022_II_15_Solution_Figure_1.png}
                    \end{figure}
                    Next, remark that $B'O_1 = DO_2$, so quadrilateral $B'O_1DO_2$ is also an isosceles trapezoid; in turn, $B'D = O_1O_2 = 15$, and similarly $A'C = 15$. Thus, Ptolemy's theorem on $A'B'CD$ yields $A'D\cdot B'C + 2\cdot 16 = 15^2$, whence $A'D = B'C = \sqrt{193}$. Let $\alpha = \angle A'B'D$. The Law of Cosines on triangle $A'B'D$ yields
 \[\cos\alpha = \frac{15^2 + 2^2 - (\sqrt{193})^2}{2\cdot 2\cdot 15} = \frac{36}{60} = \frac 35,\] and hence $\sin\alpha = \frac 45$. Thus the distance between bases $A’B’$ and $CD$ is $12$ (in fact, $\triangle A'B'D$ is a $9-12-15$ triangle with a $7-12-\sqrt{193}$ triangle removed), which implies the area of $A'B'CD$ is $\frac12\cdot 12\cdot(2+16) = 108$.

 Now let $O_1C = O_2A' = r_1$ and $O_2D = O_1B' = r_2$; the tangency of circles $\omega_1$ and $\omega_2$ implies $r_1 +  r_2 = 15$. Furthermore, angles $A'O_2D$ and $A'B'D$ are opposite angles in cyclic quadrilateral $B'A'O_2D$, which implies the measure of angle $A'O_2D$ is $180^\circ - \alpha$. Therefore, the Law of Cosines applied to triangle $\triangle A'O_2D$ yields
 \begin{align*} 193 &= r_1^2 + r_2^2 - 2r_1r_2(-\frac 35) = (r_1^2 + 2r_1r_2 + r_2^2) - \frac45r_1r_2\\ &= (r_1+r_2)^2 - \frac45 r_1r_2 = 225 - \frac45r_1r_2. \end{align*}


 Thus $r_1r_2 = 40$, and so the area of triangle $A'O_2D$ is $\frac12r_1r_2\sin\alpha = 16$.

 Thus, the area of hexagon $ABO_{1}CDO_{2}$ is $108 + 2\cdot 16 = \boxed{140}$.

 \textbf{Additional Graph for Better Understanding (Rearranging of the Pizza Slices)} This is how the reflection mentioned above is actually done. The reflection is actually a reorganization of the red and blue triangles, creating a symmetric figure with a isosceles trapezoid without changing the area. Basically you rearrange them so that each side contains a red triangle above a blue one. Then you calculate the area of the trapezoid and the two congruent triangles beside. 

 
                    \begin{figure}[H]
                        \centering
                        \includegraphics[width=0.5\textwidth]{Solution/AIME_2022_II_15_Solution_Figure_2.PNG}
                    \end{figure}
                    

 





\textbf{Solution 2}

Denote by $O$ the center of $\Omega$.Denote by $r$ the radius of $\Omega$.

 We have $O_1$, $O_2$, $A$, $B$, $C$, $D$ are all on circle $\Omega$.

 Denote $\angle O_1 O O_2 = 2 \theta$.Denote $\angle O_1 O B = \alpha$.Denote $\angle O_2 O A = \beta$.

 Because $B$ and $C$ are on circles $\omega_1$ and $\Omega$, $BC$ is a perpendicular bisector of $O_1 O$. Hence, $\angle O_1 O C = \alpha$.

 Because $A$ and $D$ are on circles $\omega_2$ and $\Omega$, $AD$ is a perpendicular bisector of $O_2 O$. Hence, $\angle O_2 O D = \beta$.

 In $\triangle O O_1 O_2$, 
 \[ O_1 O_2 = 2 r \sin \theta .  \]

 Hence, 
 \[ 2 r \sin \theta = 15 .  \]

 In $\triangle O AB$, 
 \begin{align*} AB & = 2 r \sin \frac{2 \theta - \alpha - \beta}{2} \\ & = 2 r \sin \theta \cos \frac{\alpha + \beta}{2} - 2 r \cos \theta \sin \frac{\alpha + \beta}{2} \\ & = 15 \cos \frac{\alpha + \beta}{2} - 2 r \cos \theta \sin \frac{\alpha + \beta}{2} . \end{align*}


 Hence,
 \[ 15 \cos \frac{\alpha + \beta}{2} - 2 r \cos \theta \sin \frac{\alpha + \beta}{2} = 2 . \hspace{1cm} (1) \]

 
In $\triangle O CD$, 
 \begin{align*} CD & = 2 r \sin \frac{360^\circ - 2 \theta - \alpha - \beta}{2} \\ & = 2 r \sin \left( \theta + \frac{\alpha + \beta}{2} \right) \\ & = 2 r \sin \theta \cos \frac{\alpha + \beta}{2} + 2 r \cos \theta \sin \frac{\alpha + \beta}{2} \\ & = 15 \cos \frac{\alpha + \beta}{2} + 2 r \cos \theta \sin \frac{\alpha + \beta}{2} . \end{align*}


 Hence,
 \[ 15 \cos \frac{\alpha + \beta}{2} + 2 r \cos \theta \sin \frac{\alpha + \beta}{2} = 16 . \hspace{1cm} (2) \]

 Taking $\frac{(1) + (2)}{30}$, we get $\cos \frac{\alpha + \beta}{2} = \frac{3}{5}$.Thus, $\sin \frac{\alpha + \beta}{2} = \frac{4}{5}$.

 Taking these into (1), we get $2 r \cos \theta = \frac{35}{4}$.Hence,
 \begin{align*} 2 r & = \sqrt{ \left( 2 r \sin \theta \right)^2 + \left( 2 r \cos \theta \right)^2} \\ & = \frac{5}{4} \sqrt{193} .  \end{align*}


 Hence, $\cos \theta = \frac{7}{\sqrt{193}}$.

 In $\triangle O O_1 B$,  
 \[ O_1 B = 2 r \sin \frac{\alpha}{2} .   \]

 In $\triangle O O_2 A$, by applying the law of sines, we get
 \[ O_2 A = 2 r \sin \frac{\beta}{2} .   \]

 Because circles $\omega_1$ and $\omega_2$ are externally tangent, $B$ is on circle $\omega_1$, $A$ is on circle $\omega_2$,
 \begin{align*} O_1 O_2 & = O_1 B + O_2 A \\ & = 2 r \sin \frac{\alpha}{2} + 2 r \sin \frac{\beta}{2} \\ & = 2 r \left( \sin \frac{\alpha}{2} + \sin \frac{\beta}{2} \right) .   \end{align*}


 Thus, $\sin \frac{\alpha}{2} + \sin \frac{\beta}{2} = \frac{12}{\sqrt{193}}$.

 Now, we compute $\sin \alpha$ and $\sin \beta$.

 Recall $\cos \frac{\alpha + \beta}{2} = \frac{3}{5}$ and $\sin \frac{\alpha + \beta}{2} = \frac{4}{5}$.Thus, $e^{i \frac{\alpha}{2}} e^{i \frac{\beta}{2}} = e^{i \frac{\alpha + \beta}{2}} =  \frac{3}{5} + i \frac{4}{5}$.

 We also have
 \begin{align*} \sin \frac{\alpha}{2} + \sin \frac{\beta}{2} & = \frac{1}{2i} \left( e^{i \frac{\alpha}{2}} - e^{-i \frac{\alpha}{2}}  + e^{i \frac{\beta}{2}} - e^{-i \frac{\beta}{2}}  \right) \\ & = \frac{1}{2i} \left( 1 - \frac{1}{e^{i \frac{\alpha + \beta}{2}} } \right) \left( e^{i \frac{\alpha}{2}} + e^{i \frac{\beta}{2}} \right) .   \end{align*}


 Thus,
 \begin{align*} \sin \alpha + \sin \beta  & = \frac{1}{2i} \left( e^{i \alpha} - e^{-i \alpha} + e^{i \beta} - e^{-i \beta}  \right) \\ & = \frac{1}{2i} \left( 1 - \frac{1}{e^{i \left( \alpha + \beta \right)}} \right) \left( e^{i \alpha} + e^{i \beta} \right) \\ & = \frac{1}{2i} \left( 1 - \frac{1}{e^{i \left( \alpha + \beta \right)}}  \right) \left( \left( e^{i \frac{\alpha}{2}} + e^{i \frac{\beta}{2}} \right)^2 - 2 e^{i \frac{\alpha + \beta}{2}} \right) \\ & = \frac{1}{2i} \left( 1 - \frac{1}{e^{i \left( \alpha + \beta \right)}}  \right) \left( \left(  \frac{2 i \left( \sin \frac{\alpha}{2} + \sin \frac{\beta}{2} \right)}{1 - \frac{1}{e^{i \frac{\alpha + \beta}{2}} }} \right)^2 - 2 e^{i \frac{\alpha + \beta}{2}} \right) \\ & = - \frac{1}{i} \left( e^{i \frac{\alpha + \beta}{2}} - e^{-i \frac{\alpha + \beta}{2}} \right) \left( \frac{2 \left( \sin \frac{\alpha}{2} + \sin \frac{\beta}{2} \right)^2} {e^{i \frac{\alpha + \beta}{2}} + e^{-i \frac{\alpha + \beta}{2}} - 2 } + 1 \right) \\ & = \frac{167 \cdot 8}{193 \cdot 5 } . \end{align*}


 Therefore,
 \begin{align*} {\textrm Area} \ ABO_1CDO_2 & = {\textrm Area} \ \triangle O_3 AB + {\textrm Area} \ \triangle O_3  BO_1 + {\textrm Area} \ \triangle O_3 O_1 C \\ & \quad + {\textrm Area} \ \triangle O_3 C D + {\textrm Area} \ \triangle O_3 D O_2 + {\textrm Area} \ \triangle O_3 O_2 A \\ & = \frac{1}{2} r^2 \left( \sin \left( 2 \theta - \alpha - \beta \right) + \sin \alpha + \sin \alpha + \sin \left( 360^\circ - 2 \theta - \alpha - \beta \right) + \sin \beta + \sin \beta \right) \\ & = \frac{1}{2} r^2 \left( \sin \left( 2 \theta - \alpha - \beta \right) - \sin \left( 2 \theta + \alpha + \beta \right) + 2 \sin \alpha + 2 \sin \beta  \right) \\ & = r^2 \left( - \cos 2 \theta \sin \left( \alpha + \beta \right) + \sin \alpha + \sin \beta \right) \\ & = r^2 \left( \left( 1 - 2 \cos^2 \theta \right) 2 \sin \frac{\alpha + \beta}{2} \cos \frac{\alpha + \beta}{2}  + \sin \alpha + \sin \beta \right) \\ & = \boxed{\textbf{(140) }} . \end{align*}


 





\textbf{Solution 3}


                    \begin{figure}[H]
                        \centering
                        \includegraphics[width=0.5\textwidth]{Solution/AIME_2022_II_15_Solution_Figure_3.png}
                    \end{figure}
                    Let points $A'$ and $B'$ be the reflections of $A$ and $B,$ respectively, about the perpendicular bisector of $O_1 O_2.$ 
 \[B'O_2  = BO_1 = O_1 P = O_1 C,\]
 \[A'O_1  = AO_2 = O_2 P = O_2 D.\]We establish the equality of the arcs and conclude that the corresponding chords are equal
 \[\overset{\Large\frown} {CO_1} + \overset{\Large\frown} {A'O_1} +\overset{\Large\frown} {A'B'} = \overset{\Large\frown} {B'O_2} +\overset{\Large\frown} {A'O_1} +\overset{\Large\frown} {A'B'} =\overset{\Large\frown} {B'O_2} +\overset{\Large\frown} {DO_2} +\overset{\Large\frown} {A'B'}\]
 \[\implies A'D = B'C = O_1 O_2 = 15.\]Similarly $A'C = B'D \implies \triangle A'CO_1 = \triangle B'DO_2.$

 Ptolemy's theorem on $A'CDB'$ yields 
 \[B'D \cdot A'C + A'B' \cdot CD = A'D \cdot B'C \implies\]
 \[B'D^2 + 2 \cdot 16 = 15^2  \implies B'D = A'C = \sqrt{193}.\]The area of the trapezoid $A'CDB'$ is equal to the area of an isosceles triangle with sides $A'D = B'C = 15$ and  $A'B' + CD = 18.$

 The height of this triangle is $\sqrt{15^2-9^2} = 12.$ The area of $A'CDB'$ is $108.$

 
 \[\sin \angle B'CD  = \frac{12}{15} = \frac{4}{5},\]
 \[\angle B'CD +  \angle B'O_2 D = 180^o \implies \sin \angle B'O_2 D =  \frac{4}{5}.\]

 Denote $\angle B'O_2 D = 2\alpha.$$\angle B'O_2 D > \frac{\pi}{2},$ hence $\cos \angle B'O_2 D = \cos 2\alpha = -\frac{3}{5}.$ 
 \[\tan \alpha =\frac { \sin 2 \alpha}{1+\cos 2 \alpha} = \frac {4/5}{1 - 3/5}=2.\]

 Semiperimeter of $\triangle B'O_2 D$ is           $s = \frac {15 + \sqrt{193}}{2}.$

 The distance from the vertex $O_2$ to the tangent points of the inscribed circle of the triangle $B'O_2 D$ is equal  $s – B'D = \frac{15 – \sqrt{193}}{2}.$

 The radius of the inscribed circle is $r = (s – B'D) \tan \alpha.$

 The area of triangle $B'O_2 D$ is $[B'O_2 D] = sr  = s (s – B'D) \tan \alpha = \frac {15^2 – 193}{2} = 16.$

 The hexagon $ABO_1 CDO_2$ has the same area as hexagon  $B'A'O_1 CDO_2.$ 

 The area of hexagon  $B'A'O_1 CDO_2$ is equal to the sum of the area of the trapezoid $A'CDB'$ and the areas of two equal triangles $B'O_2 D$ and $A'O_1 C,$ so the area of the hexagon $ABO_1 CDO_2$   is 
 \[108 + 16 + 16 = \boxed{140}.\]

 \textbf{vladimir.shelomovskii@gmail.com, vvsss}

 





\textbf{Solution 4}


                    \begin{figure}[H]
                        \centering
                        \includegraphics[width=0.5\textwidth]{Solution/AIME_2022_II_15_Solution_Figure_4.jpg}
                    \end{figure}
                    

 
Let circle $O_1$'s radius be $r$, then the radius of circle $O_2$ is $15-r$. Based on Brahmagupta's Formula, 

 
the hexagon's Area $= \sqrt{14(1)(16-r)(1+r)} + \sqrt{7(8)(23-r)(8+r)}$. 

 
Now we only need to find $r$.

 
Draw lines $O_1A$ and $O_1D,$ and let $X$  be the point of intersection between  $O_1D$  and  circle $O_2.$ By the SSA Congruence theorem, we have $\triangle{O_1AO_2}\cong \triangle{O_1XO_2},$ so $O_1X = O_1A.$ Thus, by the Power of A Point Theorem:

 
$(O_1A)(O_1D) = r(15+15-r) = 15r + r(15-r)$              (1).

 
Similarly,

 
$(O_2B)(O_2C) = (15-r)(15+r) = 15(15-r) + r(15-r)$       (2).

 
We can also get two other equations about these 4 segments from Ptolemy's Theorem:

 
$(O_1A)(O_2B) = 30 + r(15-r)$              (3)

 
$(O_1D)(O_2C) = 240 + r(15-r)$             (4)

 
Multiply equations (1) and (2), and equations (3) and (4) respectively, we will get a very simple and nice equation of $r$:

 
$2(15^2)r(15-r) = 7200 + 270r(15-r)$, 

 
then: 

 
$r(15-r) = 40$. 

 
This result is good enough for us to find the hexagon's area, which:

 \begin{verbatim}
	$= \sqrt{14(1)(16-r)(1+r)} + \sqrt{7(8)(23-r)(8+r)}$     $= \sqrt{14(1)(1+15-r)(1+r)} + \sqrt{7(8)(8+15-r)(8+r)}$     $= \sqrt{14(1)(1+15+40)} + \sqrt{7(8)(64+8(15)+40)}$     $= 28 + 112 = \boxed{\textbf{140}}$.
\end{verbatim}eJMaSc

 





\textbf{Solution 5 (Super bash)}


                    \begin{figure}[H]
                        \centering
                        \includegraphics[width=0.5\textwidth]{Solution/AIME_2022_II_15_Solution_Figure_5.png}
                    \end{figure}
                    

 Start by noting that the radius of $\omega_1$ and $\omega_2$ sum up to 15.

 Let the radius of $\omega_1$ be $x$, then the radius of $\omega_2$ will be $15-x$. We can find the area of the hexagon as the sum of the area of the quadrilateral $[O_1BAO_2]$ and $[O_1O_2DC]$, which is given by the Brahmagupta's Formula as: $\sqrt{(16-2)(16-15)(16-x)(1+x)} + \sqrt{(23-15)(23-16)(23-x)(8+x)}$, we only need $x$, and the condition is sufficient already. Note that the total number of diagonals and sides of the hexagon is $\frac{6 \times 5}{2} = 15$, 2 pairs of them $O_1B, O_1C$, and $O_2A, O_2D$ can all be represented by the same variable $x$, and 3 of them ($AB, O_1O_2, CD$) are given. Therefore, the total number of variables is $15-4+1-3=9$. And we have $\frac{6 \times 5}{2} = 15$ equations by doing ptolemy on the quadrilaterals inside the hexagon, we can list some of them out in hope of bashing for a solution:

 
 \[2\times 15 + x(15-x) = BO_2\times AO_1\]
 \[15\times 16 + x(15-x) = DO_1\times CO_2\]
 \[AC\times x + 2x = BC\times AO_1\]
 \[AC\times (15-x) + 16(15-x) = AD\times CO_2\]
 \[BD\times x + 16x = BC\times DO_1\]
 \[BD\times (15-x) + 2(15-x) = AD\times BO_2\]

 

 \[2CO_2 + BC\times (15-x) = AC\times BO_2\]
 \[2DO_1 + AD\times x = BD\times AO_1\]
 \[BO_2\times x + CO_2\times x = 15BC\]
 \[AD\times x + 16AO_1 = AC\times DO_1\]
 \[BC\times (15-x) + 16BO_2 = BD\times CO_2\]
 \[AO_1\times (15-x) + DO_1\times (15-x) = 15AD\]

 
 \[AO_1\times CO_2 + x(15-x) = 15AC\]
 \[BC\times AD + 2\times 16 = AC\times BD\]
 \[BO_2\times DO_1 + x(15-x) = 15BD\]

 15 equations, 9 variables, hopefully one can solve this. 

 The question is especially disgusting in the sense that there is not a lot of geometric properties or hypothesis that one can conduct. The constructions seemed hard even to USAJMO level and lacked motivation. The normal trig took way too much time and this solution took even more. The question is a straightforward and brutal one when using ptolemy, all the conditions given and inherent are used and it guarantees a solution.

 


 





\textbf{Solution 6 (Short)}

Define $r_1$ and $r_2$ as the radii of $\omega_1$ and $\omega_2$. Then, let $O_1A$ and $O_2B$ intersect at $M$ and $O_1D$ and $O_2C$ intersect at $N$. Since $AO_2 = O_2D$ and $BO_1 = O_1C$, $\angle O_1O_2=\angle O_2O_1D$ and $\angle BO_2O_1=\angle O_1O_2C$. Thus, $\triangle MO_1O_2\cong\triangle NO_1O_2$, so 
 \[\triangle MBA \sim \triangle MO_1O_2 \cong \triangle NO_1O_2 \sim \triangle NCD.\]

 Now, reflect the diagram about $\overline{O_1O_2}$ and let $a=MA$ and $b=MB$. Also, let the images of $C$ and $D$ about the reflection be $C'$ and $D'$. Notice that $C'\in\omega_1$ and $D'\in\omega_2$. Since $\angle O_1O_2B=\angle O_1O_2C$, $C'$ lies on $\overline{O_2MB}$, and similarly for $D'$ and $\overline{O_2MA}$. Now, 
 \[\text{Pow}_{\omega_1}M=MC'\cdot MB = NC\cdot MB = 8b^2=MO_1^2 - r_1^2.\] And, since $\triangle MBA \sim \triangle MO_1O_2$, $MO_1 = 15b/2$. Thus, $r_1 = \frac{\sqrt{193}b}{2}$. Similarly, $r_2 = \frac{\sqrt{193}a}{2}$. 

 To find $[AO_2DCO_1B]$, let $\theta$ be the acute angle formed by the diagonals of $ABO_1O_2$. By LoC on $\triangle BMO_1$, we find that $\cos\theta = \frac 35$. Now, \begin{align*}[AO_2DCO_1B]&=[AO_2O_1B]+[CO_1O_2D]\\&=\frac 12\sin\theta\left(a + \frac{15b}{2}\right)\left(b + \frac{15a}{2}\right) + \frac 12\sin\theta\left(8a + \frac{15b}{2}\right)\left(8b + \frac{15a}{2}\right)\\&= 27a^2 + 71ab + 27b^2.\\\end{align*}
Finally, by LoC on $\triangle MBA$, \begin{align*}27a^2 + 71ab + 27b^2 &= \frac{193}{4}(a^2 + 2ab + b^2) - \frac{85}{4}\left(a^2 + \frac{6}{5}ab + b^2\right)\\&=\frac{193}{4}\cdot\frac{4}{193}(r_1 + r_2)^2-\frac{85}{4}\cdot 4\\&= 15^2 - 85\\&=\boxed{140}.\end{align*}
 

 


 




\end{enumerate}